\documentclass[11pt]{article}
\usepackage[T1]{fontenc}
\usepackage[utf8]{inputenc}
\usepackage{lmodern}
\usepackage{amsmath,amssymb,amsthm,mathtools,mathrsfs}
\usepackage[margin=27mm]{geometry}
\usepackage{booktabs,longtable,array,graphicx}
\usepackage[hidelinks]{hyperref}
\usepackage[expansion=false]{microtype}
\setlength{\emergencystretch}{2em}
\raggedbottom
\numberwithin{equation}{section}
\newtheorem{theorem}{Theorem}[section]
\newtheorem{lemma}[theorem]{Lemma}
\newtheorem{proposition}[theorem]{Proposition}
\newtheorem{corollary}[theorem]{Corollary}
\theoremstyle{definition}
\newtheorem{definition}[theorem]{Definition}
\newtheorem{remark}[theorem]{Remark}
\newcommand{\dd}{\,\mathrm d}
\newcommand{\ii}{\mathrm i}
\newcommand{\T}{\mathbb T}
\newcommand{\R}{\mathbb R}
\newcommand{\C}{\mathbb C}
\newcommand{\Q}{\mathbb Q}
\newcommand{\Z}{\mathbb Z}
\newcommand{\eps}{\varepsilon}
\newcommand{\Pf}{\operatorname{Pf}}
\newcommand{\Res}{\operatorname{Res}}
\newcommand{\Li}{\mathcal L}
\newcommand{\dist}{\operatorname{dist}}
\newcommand{\sstar}{t_*}
\newcommand{\se}{t_\varepsilon}
\newcommand{\cO}{\mathcal O}
\newcommand{\Hpair}{\mathcal H_{\alpha\beta}}
\hypersetup{pdftitle={High-temperature bulk Ising susceptibility: a natural-boundary proof},pdfauthor={Dehao Lin}}
\title{High-temperature bulk Ising susceptibility:\\a natural-boundary proof}
\author{Dehao Lin\thanks{\raggedright School of Physics, Sun Yat-sen University.\\ORCID: \href{https://orcid.org/0009-0001-4551-8490}{0009-0001-4551-8490}. Contact: \href{mailto:lindh9@mail2.sysu.edu.cn}{\texttt{lindh9@mail2.sysu.edu.cn}}.}}
\date{17 October 2026\\Version HT-1.0\\Reader edition: 26 October 2026}
\begin{document}
\maketitle
\begin{abstract}
We prove that the complete zero-field, isotropic, infinite square-lattice
Ising bulk susceptibility, continued from the high-temperature interval,
has the unit circle as a natural boundary. The obstruction is a dense
family of prime-order Nickel singularities. At each selected point, a
finite derivative of the first odd-particle coefficient has a nonzero
inverse-distance leading term. Uniform absolute estimates show that the
entire higher odd-particle sum, including the onsite contribution, is too
small to cancel it. Signed product kernels treat both signs of the real
part of the boundary point. The analytic argument combines odd Pfaffians,
a weighted contour deformation with its Stokes current, and estimates
uniform over all particle numbers. All of these arguments are given
below, adapting the contour method of the low-temperature manuscript
without using its natural-boundary theorem.
\end{abstract}

\section{Statement, observable, and proof mechanism}\label{sec:scope}
Let $J>0$, $s=\sinh(2\beta_{\rm phys}J)$, and let
\[
 \mathcal X_H(s)=\beta_{\rm phys}^{-1}\chi_H(s)
 =\sum_{(m,n)\in\Z^2}\langle\sigma_{00}\sigma_{mn}\rangle_s
 \qquad(0<s<1).
\]
The state here is the unique zero-field infinite-volume high-temperature
Gibbs state; its magnetization vanishes. The sum runs over the entire
square lattice. Write $t=s^{-1}$ and $S(t)=t+t^{-1}$. The auxiliary
exterior variable $t$ is used only as a change of variables: no
identification of high- and low-temperature susceptibilities is made.

\begin{theorem}\label{thm:main}
The germ of $\chi_H$ inherited from $0<s<1$ has the unit circle
$|s|=1$ as a natural boundary. Equivalently, it admits no local
holomorphic continuation across any point of that circle.
\end{theorem}

The proof constructs, densely on both upper quadrants of $|t|=1$,
a point $t_*=e^{\ii\theta}$, an odd prime $p$, and
$k=(p^2-1)/2$ such that
\begin{align}
 C_p^{(k)}((1+\eps)t_*)&=4L_*\eps^{-1}+o(\eps^{-1}),
       &L_*&\ne0,\label{eq:roadmapfirst}\\
 \sum_{\substack{N>p\\N\ {\rm odd}}}
 |C_N^{(j)}((1+\eps)t_*)|&=o(\eps^{-1/2}),
       &0&\le j\le k.\label{eq:roadmaptail}
\end{align}
All lower-particle derivatives needed here are bounded. The first term's
$(k-1)$st derivative can grow logarithmically; every lower derivative
is still $o(\eps^{-1})$. These statements rule out cancellation in the
infinite sum and survive multiplication by the physical prefactor.

\paragraph{Provenance and mathematical inputs.}
The physical form-factor identity below is the classical full bulk
identity, in the precise normalization of \cite[eqs.~(1.9), (1.11),
(3.13), (3.17)--(3.20)]{Orrick2001} and \cite[eqs.~(3), (5)--(9)]{Boukraa2008}.
The contour partition and tail estimates adapt the method of manuscript
v0.1-rc6, at commit
\texttt{2cf997226db314ad52fa40a2ecda5439542606bf} \cite{Snapshot};
the adapted arguments are proved here. The physical identity is the
external statistical-mechanical input. The remaining argument derives
the analytic obstruction from the odd-particle integrals. In particular,
the fixed-order gradient method of \cite{TW2014} is established at the
selected odd orders, including the additional negative-real-part endpoint
case. The proof does not use the low-temperature natural-boundary theorem.

\paragraph{How to read the proof.}
The geometric picture is a boundary point at which the first admissible
particle contribution becomes singular, while the full remaining series
stays smaller after the same finite number of derivatives. Sections
\ref{sec:setup}--\ref{sec:first} identify the complete lattice observable,
select these points, and calculate the first singularity. Section
\ref{sec:signs} reduces both boundary halves to one signed family.
Sections \ref{sec:tailreduction}--\ref{sec:sectors} supply the uniform
estimates needed for the infinite sum; Section \ref{sec:conclusion}
assembles them, and Section 11 completes the noncancellation argument.
Appendix \ref{app:lie} gives the differentiation and coarea calculus used
in the sector estimates. Appendix \ref{app:constants} records the order
in which all constants are fixed.

A first reading can follow the main theorem, the first-coefficient
Theorem~\ref{thm:first}, the tail Theorem~\ref{thm:tail}, and their
assembly in Section 11. The intervening sections provide the complete
proof of the tail theorem, including the contour current and the onsite
series.

\tableofcontents

\section{The full odd-particle representation}\label{sec:setup}

The angular formula identifies the physical observable. The double
contour formula is the form used for estimates. The bridge between them
is the exact full-site identity in Lemma~\ref{lem:residue}: it keeps the
origin and all reflected axes and quadrants. Lemma~\ref{lem:schur} then
justifies differentiation of the complete series on compact subsets of
the exterior domain.
For real $t>1$, put
\[
 W(y;t)=S(t)-\tfrac12(y+y^{-1}),\qquad
 z=W-\sqrt{W^2-1}=e^{-\gamma},\qquad
 R(y;t)=\frac{2z^2}{1-z^2}.
\]
On the angular unit circle, the root is specified by $\gamma>0$.
Elsewhere it is its continued interior root. For
$\varphi_N=-\sum_{i<N}\varphi_i\pmod{2\pi}$, set
\begin{equation}
 C_N(t)=\frac1{N!}\int\prod_{i<N}\frac{\dd\varphi_i}{2\pi}
 \prod_{i=1}^N\frac1{\sinh\gamma_i}
 \frac{1+Z}{1-Z}
 \prod_{i<j}\left[
 \frac{\sin((\varphi_i-\varphi_j)/2)}
      {\sinh((\gamma_i+\gamma_j)/2)}\right]^2,
 \quad Z=\prod_i e^{-\gamma_i}.
 \label{eq:angular}
\end{equation}
Here $\cosh\gamma_i=S(t)-\cos\varphi_i$, and every independent
angle is integrated over a full period. The physical identity is
\begin{equation}
 \mathcal X_H(s)=P_H(s)\sum_{\substack{N\ge1\\N\ {\rm odd}}}C_N(1/s),
 \qquad P_H(s)=\frac{(1-s^4)^{1/4}}s.
 \label{eq:bulk}
\end{equation}
The fourth root is fixed by positivity on $0<s<1$. The sum in
\eqref{eq:bulk} includes every positive odd particle number and every
lattice site. The one-particle value provides a normalization check:
$C_1=1/(S-2)=s/(1-s)^2$, fixing the normalization.

Define $D(x,y;t)=S(t)-(x+x^{-1}+y+y^{-1})/2$. On sufficiently
large interior circles $|x_i|=|y_i|=r<1$, define
\begin{align}
 T_N(t)&=\frac1{N!}\oint\cdots\oint
 \frac{X^{-1}+Y^{-1}}{(1-X)(1-Y)}
 \prod_{i<j}\frac{x_i-x_j}{1-x_ix_j}\frac{y_i-y_j}{1-y_iy_j}
 \prod_i\frac{\dd x_i\,\dd y_i}{(2\pi\ii)^2D(x_i,y_i;t)},
 \label{eq:double}\\
 X&=\prod_i x_i,\qquad Y=\prod_i y_i.\nonumber
\end{align}
Let $f_N$ be the same integral with its global rational factor
replaced by $1/(XY)$.

\begin{lemma}[Residue reduction and the full-site identity]\label{lem:residue}
For every positive integer $N$, the residue-reduced expression is
\begin{equation}
 T_N=\frac1{N!}\oint\cdots\oint
 \frac{Z^{-1}+Y^{-1}}{(1-Z)(1-Y)}
 \prod_{i<j}P_{ij}\prod_i R(y_i;t)\frac{\dd y_i}{2\pi\ii},
 \label{eq:reduced}
\end{equation}
where
\begin{equation}
 P_{ij}=\frac{z_i-z_j}{1-z_iz_j}\frac{y_i-y_j}{1-y_iy_j}
 =-\frac{(y_i-y_j)^2z_iz_j}{y_iy_j(1-z_iz_j)^2}.
 \label{eq:pair}
\end{equation}
The onsite coefficient $f_N$ has the reduced global factor $1/(ZY)$,
and
\begin{equation}
 C_N=f_N+2T_N.\label{eq:onsiteidentity}
\end{equation}
No parity assumption is used in these identities.
\end{lemma}
The proof first establishes the contour identity and then attaches the
full lattice sum; both steps apply independently of parity.
\begin{proof}
\noindent\textbf{Residues and pair cancellation.} Choose $r$ with $|z(y;t)|<r$ on the $y$ circle. The only enclosed
dispersion root in each $x_i$ is $z_i$, with residue
$2z_i^2/(1-z_i^2)$. At $x_i=0$, $D^{-1}=-2x_i+O(x_i^2)$ cancels
the apparent numerator pole. Pair and product poles involving the
other interior variables remain outside the $x_i$ circle after each
successive residue. Subtracting the two dispersion equations gives
\[
 (z_i-z_j)(1-z_iz_j)
 =-\frac{z_iz_j}{y_iy_j}(y_i-y_j)(1-y_iy_j),
\]
which proves the canceled pair expression.

\par\smallskip\noindent\textbf{Reconstructing the lattice sum.} On $y_i=e^{\ii\varphi_i}$, $P_{ij}$ is precisely the squared pair
in \eqref{eq:angular}, and
\[
 R(y_i;t)\frac{\dd y_i}{2\pi\ii}
 =y_i z_i\frac{\dd\varphi_i}{2\pi\sinh\gamma_i}.
\]
Let $q_N(m,n)$ be the resulting $N$-angle integral with multiplier
$Z^mY^n$, for $m\ge0$, $n\in\Z$. Angle reversal gives
$q_N(m,n)=q_N(m,-n)$. Its double-contour lift gives
$q_N(m,n)=q_N(n,m)$ for nonnegative $m,n$: exchange the two sets
of variables, using Fubini with the induced orientation in each
coordinate. Any sign of the exchanged top form is canceled by the
same sign of the exchanged product orientation.

For fixed real $t>1$, a small common complex angle strip has
$\Re\gamma\ge c>0$ and bounded complete pairs. A simultaneous
angle shift therefore bounds $q_N(m,n)$ by a constant depending on
$N,t$ times $e^{-N(c m+\eta|n|)}$. The sums over lattice
coordinates are absolutely convergent. Their horizontal Fourier
sum enforces $\sum\varphi_i=0\pmod{2\pi}$, while their vertical
sum is $(1+Z)/(1-Z)$. This gives \eqref{eq:angular}.
The reduced scalar multiplier in $T_N$ is
\[
 \frac{Y+Z}{(1-Y)(1-Z)}
 =\sum_{m\ge1,n\ge0}Y^mZ^n+
   \sum_{m\ge0,n\ge1}Y^mZ^n.
\]
It counts each positive axis once and each open positive quadrant
twice. Thus $f_N+2T_N$ counts the origin once, each reflected axis
twice, and each reflected open quadrant four times, exactly as the
full lattice sum. This proves \eqref{eq:onsiteidentity} on the
real interval and hence for its holomorphic continuation.
\end{proof}

\begin{lemma}[Odd Pfaffians and normal convergence]\label{lem:schur}
For even $m$,
$\Pf[(v_i-v_j)/(1-v_iv_j)]_{i,j\le m}
=\prod_{i<j}(v_i-v_j)/(1-v_iv_j)$.
For odd $N$, append $v_{N+1}=-1$ and let $\widehat A(v)$ be the
resulting even matrix, whose last upper column is identically one.
Then
\[
 \Pf\widehat A(v)=\prod_{1\le i<j\le N}\frac{v_i-v_j}{1-v_iv_j}.
\]
The series of $T_N,f_N,C_N$, and each of their fixed derivatives,
converge locally normally for $|t|>1$.
\end{lemma}
\begin{proof}
\noindent\textbf{Pfaffian identity.} For the even identity, set $x_i=(1+v_i)/(1-v_i)$ and induct on
$m$. As a rational function of $x_m$, both sides have the same
simple residues at $x_m=-x_i$: expansion of the Pfaffian and
evaluation of the product give the sign $(-1)^{i+1}$ and the
factor $2x_i$ times the smaller identity. Their difference is a
bounded polynomial and vanishes at $x_m=x_{m-1}$, hence is zero.
The odd formula follows by the indicated substitution, with
removable cases obtained by continuity.

\par\smallskip\noindent\textbf{Normal convergence and derivatives.} On a compact exterior neighborhood and a slightly larger one choose
$r<1$ and $0<m_0\le |z|\le q<r$ uniformly. Such choices follow
from the interior-root equation, compactness, and then a small
decrease of the unit $y$ radius. Every one-body factor and every
Schur entry is bounded. The inverse products cost at most $C^N$,
and the global product denominators have fixed positive gaps.
Hadamard's inequality and $(\Pf A)^2=\det A$ bound an even
matrix of size $m$ with entries at most $C$ by
$C^{m/2}(m-1)^{m/4}$. Applied to both augmented matrices, this gives
\[
 |T_N|+|f_N|+|C_N|
 \le C_0 C_1^N\frac{N^{(N+1)/2}}{N!}
 \le C_0\sqrt N\left(\frac{eC_1}{\sqrt N}\right)^N.
\]
This is summable. Cauchy's inequality on the larger compact
neighborhood proves the assertion for every fixed derivative.
\end{proof}

\section{A dense first-odd-order family}\label{sec:arithmetic}

The selected points must meet three conditions: a known first odd order,
uniqueness of the resonant cosine pair at that order, and exclusion of
endpoint resonances. The prime construction establishes these conditions
together. The final resultant estimate converts endpoint nonresonance
into the quantitative gap used in the tail proof.
Let $p\ge11$ be prime, and choose distinct
$a,b\in\{1,\ldots,(p-1)/2\}$ such that
\[
 \alpha=2\pi a/p,\quad\beta=2\pi b/p
\]
both belong to $(0,\pi/2)$ or both belong to $(\pi/2,\pi)$. Put
\[
 S_* =\cos\alpha+\cos\beta,
 \quad \theta=\arccos(S_*/2),\quad t_*=e^{\ii\theta},
 \quad N_0=p,\quad k=\frac{p^2-1}{2}.
\]
In particular $S_*\ne0$ and $k$ is divisible by four.

\begin{theorem}\label{thm:prime}
These points are dense in both open upper quadrants. Their first
odd Nickel order is $p$, and at order $p$ the unordered cosine
pair is uniquely $\{\cos\alpha,\cos\beta\}$. Neither of the
possible endpoint configurations is resonant at any particle number.
\end{theorem}
\begin{proof}
\noindent\textbf{First order.} Write $\zeta=e^{2\pi\ii/p}$, $u_j=\zeta^j+\zeta^{-j}$,
$d=(p-1)/2$, $E=\Q(\zeta+\zeta^{-1})$, and $v=u_a+u_b=2S_*$.
The trace of every nonzero-index $u_j$ is $-1$, so
$\operatorname{Tr}_E v=-2$. If the algebraic integer $v$ were
rational it would be an integer, whereas the trace would force
$v=-2/d$, impossible for $d\ge5$.

If $p\nmid n$, the cyclotomic fields $\Q(\zeta_p)$ and
$\Q(\zeta_n)$ have trivial intersection. Indeed their compositum
is $\Q(\zeta_{pn})$ of degree $\varphi(p)\varphi(n)$, the product
of their degrees. A Nickel representation
$S_* =\cos(2\pi l/n)+\cos(2\pi m/n)$ would put $S_*$ in their
intersection, a contradiction. Since order $p$ realizes our pair,
it is the first odd order.

\par\smallskip\noindent\textbf{Uniqueness and density.} The only rational relation between $1,u_1,\ldots,u_d$ is a
multiple of $1+\sum_j u_j=0$. To see this, lift the relation to
a polynomial of degree at most $p-1$ and use the minimal
polynomial $1+X+\cdots+X^{p-1}$. An equality of two cosine
pairs has at most four nonconstant positions. Because $d\ge5$,
one zero coefficient forces the multiple to vanish, proving uniqueness.
Approximation of any angle in either open half-interval by two
distinct neighboring prime-grid angles proves density.

\par\smallskip\noindent\textbf{Endpoint exclusion.} For $S_*>0$, an endpoint resonance would require a root of unity
$\xi$ with $\xi+\xi^{-1}=2(S_*-1)=v-2$. Its trace is
$-2-2d<-2d$, whereas every conjugate of a root-of-unity trace lies
in $[-2,2]$. For $S_*<0$ the corresponding equation is
$\xi+\xi^{-1}=2(S_*+1)=v+2$. Here
\[
 \operatorname{Tr}_E(v^2)=4d-6=2p-8.
\]
This follows from $u_a^2=u_{2a}+2$ and
$u_a u_b=u_{a+b}+u_{a-b}$; all the indicated indices are nonzero
modulo $p$. Were all conjugates of $v+2$ in $[-2,2]$, all
conjugates of $v$ would be in $[-4,0]$ and would satisfy
$v^2\le-4v$. Taking traces would imply $4d-6\le8$, contrary
to $d\ge5$. This excludes the negative-side endpoint as well.
\end{proof}

\begin{lemma}[Exponential nonresonance]\label{lem:resultant}
Every fixed algebraic unit-modulus number $\xi$ that is not a root
of unity obeys $|1-\xi^N|\ge e^{-AN}$ for some $A>0$ and all
$N\ge1$. After increasing $A$, both $|1-\xi^N|$ and
$|1+\xi^N|$ obey this bound.
\end{lemma}
\begin{proof}
Let $P(X)=a_d\prod_\nu(X-\xi_\nu)$ be its primitive irreducible
integer polynomial, with $\xi_1=\xi$. Its nonzero integer
resultant with $X^N-1$ has modulus at least one, so
\[
 1\le |a_d|^N|1-\xi^N|\,2^{d-1}
        \prod_{\nu>1}\max(1,|\xi_\nu|)^N.
\]
This proves the first bound. Use
$|1+\xi^N|\ge |1-\xi^{2N}|/2$ for the second.
\end{proof}

\section{First singularity and bounded complements}\label{sec:first}

We first isolate the only configurations that can obstruct bounded
fixed-order differentiation. We then calculate their contribution with
the mean integration performed before the shape estimate. This order
retains the residue normalization and produces the nonzero coefficient
needed in the final infinite-sum argument.
In this section $t_*$ can lie in either upper quadrant. Set
$t_\eps=(1+\eps)t_*$ and keep $r=e^{-c_0\eps}$ fixed during
every $t$ derivative. A sufficiently small fixed $c_0>0$ gives
$\Im W\ge c\eps$ on a $c'\eps$ parameter disk. The root lies
in the lower unit half-disk: writing $\gamma=u+\ii v$ gives
$\Im W=\sinh u\sin v\le\sinh u$, hence $u\ge c''\eps$.
Decreasing $c_0$ also gives $|z|<r$, as required for the residues.

\begin{lemma}[Fixed-particle regularity]\label{lem:complement}
All derivatives of every fixed $T_n,f_n,C_n$ with odd $n<p$
are bounded on this radial approach. At $n=p$, the complement
of the two ordered lower double-contour configurations
$(x_j,y_j)=(e^{-\ii\beta},e^{-\ii\alpha})$ and its exchange
has bounded derivatives of every fixed order. Every fixed
derivative of $f_n$ is bounded at these points, for every fixed $n$.
\end{lemma}
\begin{proof}
The active denominator-gradient vectors are the positive all-ones
vectors for the two global products, the positive reciprocal-pair
vectors, and the dispersion vectors with entries
$(\Im x_j,\Im y_j)$. The sign follows from $\Im S(t_\eps)>0$.
Suppose a nontrivial nonnegative zero combination exists. A pair
vector cannot occur: cancellation of its two positive entries would
require two strictly negative imaginary parts, while reciprocal
unit-circle coordinates have opposite imaginary parts.

If both product vectors occur, all $x_j,y_j$ lie in the lower
semicircle and their imaginary-part ratios agree. Along
$\cos A+\cos B=S_*\ne0$,
\[
 \frac{\dd}{\dd A}\frac{\sin A}{\sin B}
 =\frac{S_*(1-\cos A\cos B)}{\sin^3 B}
\]
has a fixed strict sign. Thus all the coordinate pairs are equal,
and the two product equations make both common coordinates $n$th
roots of unity. Theorem~\ref{thm:prime} excludes them below $p$
and leaves just the stated two configurations at $p$.
If only one product vector occurs, the other coordinates have zero
imaginary part and are all $+1$ for $S_*>0$ or all $-1$ for
$S_*<0$. The remaining common coordinate is a root of unity,
excluded by the endpoint argument in Theorem~\ref{thm:prime}.
The latter case is exactly the additional odd-order possibility
left in \cite[footnote~5]{TW2014}; no even-order theorem is being
applied to it. With neither product vector, a nonzero dispersion
vector cannot cancel, since $S_*\notin\{-2,0,2\}$.

Cover the allowed compact torus support by finitely many fixed
smooth charts. On each, a real direction separates all active
vectors from zero by a common positive margin, persisting on the
nearby damped contours. Represent the $\ell$ active denominators
as positive-half-line exponential integrals, and write
$R=\sum_{i=1}^{\ell}\xi_i$. The exponent $g$ has
$\Re g\le0$, $|e\cdot\nabla g|\ge cR$, and bounded fixed
angular jets of size $O(R)$. A $j$th parameter derivative costs
at most $R^j$. On $R\ge1$, integrate by parts
$q=\ell+j+1$ times with $(e\cdot\nabla)/(e\cdot\nabla g)$.
Including all derivatives of coefficients and cutoffs, the
remaining radial integral is bounded by
$C\int_1^\infty R^{\ell-1+j-q}\dd R=C\int_1^\infty R^{-2}\dd R$.
Small $R$ is directly bounded. Omitting both product factors gives
the assertion for $f_n$. Constants may depend on the fixed $n,j,t_*$.
\end{proof}

\begin{theorem}[The complete first coefficient]\label{thm:first}
Let
\[
 \Hpair=S_*(1-\cos\alpha\cos\beta)\ne0,
 \qquad P_p=\prod_{m=1}^{p-1}\frac{m!}{2^m},
\]
and define
\begin{equation}
 B_* =\frac{P_p(p\sin\theta)^{k-1}}
 {\sqrt p\,\pi^{(p-1)/2}|\Hpair|^k}>0,
 \qquad L_*=-\ii e^{-\ii k\theta}B_*.
 \label{eq:amplitudeclosed}
\end{equation}
Then
\[
 T_p^{(k)}(t_\eps)=2L_*\eps^{-1}+o(\eps^{-1}),\qquad
 C_p^{(k)}(t_\eps)=4L_*\eps^{-1}+o(\eps^{-1}).
\]
Derivatives through order $k-2$ are bounded, and those of order
$k-1$ are $O(1+|\log\eps|)$.
\end{theorem}
\begin{proof}
\noindent\textbf{1. Localize the two configurations.} For the lower $y$ chart centered at $e^{-\ii\alpha}$ write
$y_j=r e^{\ii(-\alpha+u_j)}$, $v=\sum u_j$,
$u_j=v/p+\eta_j$, $\sum\eta_j=0$. The Jacobian from
$(v,\eta_1,\ldots,\eta_{p-1})$ to $u$ is one. Choose fixed
smooth $y$-only weights $w_\alpha=\eta(v)\chi(\eta_1,\ldots)$,
and $w_\beta$ on the exchanged chart, equal to one near the two
configurations and with disjoint small supports. The exact
double-integral partition leaves a complement controlled by
Lemma~\ref{lem:complement}. All $x$ integrations stay unweighted,
so the residues in Lemma~\ref{lem:residue} apply to each chart.

\par\smallskip\noindent\textbf{2. Integrate the mean variable.} On a fixed hard mean segment (whose difference from the smooth
mean cutoff has separated denominators), put $z_j=e^{-\ii\phi_j}$.
Implicit differentiation gives
\[
 b=\phi_u=\frac{\sin\alpha}{\sin\beta}>0,\quad
 d=-\tfrac12\phi_{uu}=\frac{\Hpair}{2\sin^3\beta}\ne0,
 \quad A_0=\frac{2\sin\theta-c_0\sin\alpha}{\sin\beta}.
\]
Thus
$\phi=\beta-\ii A_0\eps+bu-du^2+
O(\eps^2+\eps|u|+|u|^3)$. Choose $c_0$ with $A_0>c_0$.
Move the hard mean segment downward on a small fixed rectangle.
The $Y=1$ pole is $v=-\ii pc_0\eps$; the $Z$ pole stays above
the real mean line. Downward mean motion increases its damping,
independently of the sign of $d$. On the vertical sides the phase
gap is $b\delta+O(|d|h^2+h_0^2)$; choose the shape radius $h$
and rectangle depth $h_0$ small relative to $\delta$. The bottom
has fixed positive damping. All three displaced sides and every
fixed parameter derivative are bounded. The clockwise closure and
$\partial_v(1-Y)=-\ii$ give the residue multiplier $+2\pi$.

\par\smallskip\noindent\textbf{3. Estimate the constrained shape integral.} At this residue the angles are exactly $y_j=e^{\ii(-\alpha+\eta_j)}$.
The reduced integrand becomes the corresponding constrained
angular chart in \eqref{eq:angular}: its multiplier at $Y=1$ is
$(1+Z)/(1-Z)$, and $2\pi$ removes one angular normalization.
Let $q=\sum\eta_j^2$. On the fixed shape ball,
\begin{equation}
 1-Z=Q\eps-\ii d q+
 O(\eps^2+\eps|\eta|+|\eta|^3),\quad
 Q=\frac{2p\sin\theta}{\sin\beta}>0,
 \quad D_t=Qe^{-\ii\theta}.
 \label{eq:localdenom}
\end{equation}
Moreover $|1-Z|\ge c(\eps+q)$. Indeed modulus damping is at
least $c\eps$, the constrained linear shape term vanishes by
permutation symmetry and $\sum\eta_j=0$, the phase at zero shape
differs from its limiting multiple of $2\pi$ by $O(\eps^2)$,
and the phase Hessian is
strictly definite with the sign of $-\Hpair$ on a sufficiently
small regular chart. Its integral remainder bounds the phase by
$cq$ outside $q=O(\eps^2)$; damping covers that remaining region.

The leading constrained numerator has phase
$(-\ii)^p(-1)^{p(p-1)/2}=-\ii$ and density
\[
 \frac{-2\ii}{p!(2\pi)^{p-1}2^{p(p-1)}\sin^{p^2}\beta}
 \frac{\Delta(\eta)^2\dd\eta_1\cdots\dd\eta_{p-1}}
 {Q\eps-\ii d\sum\eta_j^2}.
\]
After $k$ derivatives, scale $\eta=\sqrt\eps\,u$. Since
$k=(p^2-1)/2$, the scale is $\eps^{-1}$, and the radial
majorant has large-radius behavior $R^{-3}$. More explicitly,
write the exact analytic numerator as $\Delta(\eta)^2A(t,\eta)$.
Parameter derivatives preserve this Vandermonde. The differentiated
density is a finite sum with denominator powers at most $k+1$;
only the term $(-1)^k k!D_t^k\Delta^2A/(1-Z)^{k+1}$ contributes
to the limit after multiplication by $\eps$. Its rescaling is
dominated by $C\Delta(u)^2/(1+|u|^2)^{k+1}$.
For every lower pole power, the radial bound
\[
 \eps\int_0^h\frac{r^{2k-1}}{(\eps+r^2)^{m+1}}\dd r
 \longrightarrow0\qquad(0\le m<k)
\]
follows by splitting at $r=\sqrt\eps$; at $m=k-1$ the cost is
$O(\eps|\log\eps|)$. This justifies the differentiated limit,
without interchanging an undominated mean/shape absolute integral.

\par\smallskip\noindent\textbf{4. Evaluate the coefficient and add the charts.} Let $\mathscr C_p=p^{-1/2}\int_{\sum\omega_i=0,\,|\omega|=1}
\Delta(\omega)^2\dd\sigma(\omega)>0$. Polar coordinates give
\[
 \int_{\R^{p-1}}\frac{\Delta(u)^2\dd u}
 {(Q-\ii d\sum u_i^2)^{k+1}}
 =\frac{\mathscr C_p}{2}Q^{-1}(-\ii d)^{-k}B(k,1)
 =\frac{\mathscr C_p}{2kQd^k}.
\]
For an elementary evaluation set $c=-\ii d$ and $x=|u|^2$ in
the radial integral. The exact identity
\[
 \frac{\dd}{\dd x}\left(\frac{cx}{Q+cx}\right)^k
 =kQc^k\frac{x^{k-1}}{(Q+cx)^{k+1}}
\]
has endpoint values zero and one, and $Q+cx\ne0$ for $x\ge0$.
It proves the displayed integral without any continuation of
a noninteger power. Since $k\equiv0\pmod4$,
$d^k=|d|^k>0$ for either sign of $S_*$. The ordinary Gaussian
Vandermonde integral is
\[
 \int_{\R^p}e^{-|x|^2/2}\Delta(x)^2\dd x
 =(2\pi)^{p/2}\prod_{m=1}^p m!.
\]
For completeness, expand the two Vandermonde determinants in monic
Hermite polynomials. Orthogonality gives $p!$ times the product of
their squared norms $\sqrt{2\pi}\,m!$, $0\le m<p$, obtained by
integration by parts from the Rodrigues formula. Separate the unit
mean coordinate and polar coordinates on its orthogonal complement
to obtain
\[
 \mathscr C_p=
 \frac{(2\pi)^{(p-1)/2}\prod_{m=1}^p m!}
 {\sqrt p\,2^{k-1}(k-1)!}.
\]
Substitution cancels every power of $\sin\beta$ and yields exactly
$L_*$ in \eqref{eq:amplitudeclosed}. Exchanging $\alpha,\beta$
leaves this expression unchanged, so the two ordered charts add.
The fixed onsite coefficient is smooth by Lemma~\ref{lem:complement},
and \eqref{eq:onsiteidentity} gives the factor four for $C_p$.
Equivalently the full angular representation has four equal charts
at $\pm\alpha,\pm\beta$. Finally the same radial bound at derivative
order $k-1$ is logarithmic, and through order $k-2$ is integrable
at $\eps=0$. The displaced sides and complement are already bounded.
\end{proof}

\section{Signed kernels and the two boundary halves}\label{sec:signs}

For odd particle number, reflection changes the sign of a product
kernel. Introducing a sign parameter makes this change explicit and
allows the estimates on the two upper quadrants to share one contour
geometry. The sign is retained in every global denominator below.
For $\sigma\in\{1,-1\}$, define $T_N^\sigma$ by replacing
the global factor in \eqref{eq:reduced} with
\[
 \frac{Z^{-1}+Y^{-1}}{(1-\sigma Z)(1-\sigma Y)}.
\]
Then $T_N=T_N^1$. Simultaneously replacing
$(t,y,z)$ by $(-t,-y,-z)$ leaves $P_{ij},R$ unchanged;
the numerator and the product of $N$ differentials each acquire
$(-1)^N$. Consequently
\begin{equation}
 T_N^\sigma(-t)=T_N^{\sigma(-1)^N}(t),
 \qquad f_N(-t)=f_N(t)\,(-1)^N.
 \label{eq:signidentity}
\end{equation}
Conjugation also holds, by the real germ and analytic continuation.
Thus a quadrant-II estimate for odd $N$ reduces, by negation and
conjugation, to a quadrant-I estimate with $\sigma=-1$. This identity concerns the auxiliary integrals; it leaves the physical
observable $\mathcal X_H$ with its original high-temperature germ.

In the tail proof below, fix a point with $0<S_*<2$, and require
that $\xi=e^{-\ii\theta_B}$, $\cos\theta_B=S_*-1$, be algebraic
and not a root of unity. Both the original positive family and the
reflection of the negative family satisfy this by
Theorem~\ref{thm:prime}. Increase $A$ in Lemma~\ref{lem:resultant}
so $|1-\sigma\xi^N|\ge e^{-AN}$ for either sign and all $N$.
Throughout the following tail sections $T_N$ denotes the fixed
signed integral $T_N^\sigma$; the point, sign, odd $N_0=p$, and
$k=(p^2-1)/2$ are held fixed. All absolute bounds are uniform in
the particle windows stated there. The onsite estimates will not
require any resummation identity for $\sum f_N$.

\section{Exact contour decomposition and quantitative requirements}
\label{sec:tailreduction}

The deformation separates three tasks: a final contour with a selected
upper index, the original contour on the lower support, and the current
created by the weight. The exact identity below keeps all three. Its
four target estimates organize the remainder of the tail proof.
Fix the signed integral and positive-$S_*$ geometry of Section~\ref{sec:signs}, with odd $N_0=p$ and $k=(p^2-1)/2$.
Set \(H=\log(1/\eps)\) and choose
\begin{equation}
 0<\alpha_0<\frac1{2(k+2)},\qquad \lambda_*=\eps^{\alpha_0}.
 \label{eq:alphachoice}
\end{equation}
Neither \(N_0\) nor \(k\) varies with \(\eps\). Throughout the tail proof, $t_\eps=(1+\eps)t_*$ and $r=e^{-c_0\eps}$; the radius is held fixed under each parameter derivative.
Let $\Omega_{N,t}$ be the signed reduced holomorphic form of Section~\ref{sec:signs}, including its $1/N!$.
Choose smooth periodic functions \(p,m,a\) with values in \([0,1]\)
as follows, in angular representatives \(-\pi<\theta\le\pi\).
The closed support of \(m\) is contained in a small neighborhood of
\(-\theta_B\) within \((-\pi,0)\); \(m=1\) on an open neighborhood
of the closed branch-cutoff supports in Lemma~\ref{lem:partition}.
Thus it is one on the supports of all their finite jets as well.
Choose a small fixed endpoint width \(\delta>0\) and take
\[
 a=0\ \hbox{off }[\delta,\pi-\delta],\qquad
 a=1\ \hbox{on }[2\delta,\pi-2\delta].
\]
Make \(\delta\) small enough that the reflection of
\(\operatorname{supp}m\) lies in the interior of the latter interval.
Take \(p\) with closed support in \((0,\pi)\), equal to one on
an open neighborhood of \([\delta,\pi-\delta]\).
Smooth step functions give all these choices simultaneously.
Their widths and the deformation strength are fixed in the order
of Appendix~\ref{app:constants}, before \(\eps\to0\).
The supports of \(p,m\) are disjoint, the upper inverse branch is
inside \(a=p=1\), and \(\operatorname{supp}a\) and
\(\operatorname{supp}a'\) have positive distance from both endpoints.
In particular a named upper angle has \(\sin\theta\ge\sin\delta>0\).

Define, with fixed deformation strength \(\tau>0\),
\[
 P=\sum_i p(\theta_i),\quad
 h_i=-2\tau p(\theta_i)+\frac{\tau P}{2N}m(\theta_i),\quad
 y_i(\lambda,\theta)=r e^{\ii\theta_i+\lambda h_i},
 \qquad \psi(\theta)=\prod_i(1-a(\theta_i)).
\]
The letter \(P\) here is an occupancy sum, not the pair factor \(P_{ij}\).
Write \(U_{\lambda,t}=\Phi_\lambda^*\Omega_{N,t}\).

\begin{proposition}[Exact weighted retraction]\label{prop:selector}
For exterior \(t\) in its admissible neighborhood,
\begin{align}
 T_N&=F_N+K_N+S_N,\label{eq:FKS}\\
 F_N&=\int(1-\psi)U_{1,t},\qquad K_N=\int\psi U_{0,t},\nonumber\\
 S_N&=-2\ii\tau\int_0^1\sum_q\int
                         (\partial_{\theta_q}\psi)U_{\lambda,t}\dd\lambda.
 \label{eq:current}
\end{align}
The current has one named index \(q\) in \(\operatorname{supp}a'\),
with \(p_q=1,\ p'_q=m_q=m'_q=0\); all other indices belong to
the enlarged lower support \(K=\operatorname{supp}(1-a)\).
\end{proposition}
\begin{proof}
We have \(\sum_i h_i\le-3\tau P/2\), so \(|Y|<1\).
An inward radial motion on an upper angle and an outward radial motion
on a lower angle both increase \(\Im W\):
\(\partial_{\log|y|}\Im W=-\cosh(\log|y|)\sin\theta\).
Thus all \(z_i\)'s stay on the interior sheet. Neither \(1-Z\)
nor a denominator \(1-z_iz_j\) vanishes. The canceled pair expression
has no remaining \(1-y_iy_j\) pole. This proves legality for the full
closed homotopy.

Normalize the angular Jacobian row \(i\) by \(y_i\). Its matrix is
\[
 M_{ij}=d_i\delta_{ij}+\frac{\lambda\tau}{2N}m_i p'_j,\qquad
 d_i=\ii+\lambda\left(-2\tau p'_i+\frac{\tau P}{2N}m'_i\right).
\]
Because \(m_ip'_i=0\), \(\det M=\prod_i d_i\); all factors are
nonzero and its absolute value is bounded by \(C^N\).
At a named \(q\), both the row and column of \(M\) are zero
off the diagonal and \(d_q=\ii\). Replacing angular column \(q\)
by the homotopy velocity multiplies the determinant by
\((-2\tau)/\ii=2\ii\tau\).
Cartan's formula on the closed torus gives
\[
 \partial_\lambda\int\psi U_{\lambda,t}
 =-2\ii\tau\sum_q\int(\partial_{\theta_q}\psi)U_{\lambda,t}.
\]
The unweighted integral is homotopy invariant. Integrating this identity
and separating the selected part proves \eqref{eq:FKS}.
\end{proof}

The current measures the change of the weighted integral during the
homotopy; it is part of the exact decomposition. All \(t\)-derivatives
are taken on the entire fixed interval
\([0,1]\) before splitting the derivative density at \(\lambda_*\).
Thus no derivative of that split endpoint enters.

For every odd $N>N_0$, denote the already differentiated,
angularly integrated current by
\begin{gather*}
 \mathcal J_{N,j}(\lambda)
   =-2\ii\tau\sum_q\int(\partial_{\theta_q}\psi)\,
              \partial_t^j U_{\lambda,t}\big|_{t=t_\eps},\\
 B_{N,j}=|K_N^{(j)}(t_\eps)|
       +\left|\int_0^{\lambda_*}\mathcal J_{N,j}(\lambda)\dd\lambda\right|.
\end{gather*}
The quantitative requirements are
\begin{align}
 \sum_N|F_N^{(j)}|
   &\le \exp(C_j\log^2(H+2)),\label{eq:reqF}\\
 \sum_N\left|\int_{\lambda_*}^1
               \mathcal J_{N,j}(\lambda)\dd\lambda\right|
   &\le C_j\eps^{-\alpha_0(j+2)},\label{eq:reqLS}\\
 B_{N,j}
   &\le C_j^N\eps^{-j-2}e^{-\kappa N^2}
             &&(N\ge D\sqrt H),\label{eq:reqhigh}\\
 B_{N,j}
   &\le \exp(C_jN\log(N+1))(H+N)^{C_j}
             &&(N_0+2\le N<D\sqrt H).\label{eq:reqinter}
\end{align}
The sums are over all odd tail orders \(N>N_0\), without an upper particle-number cutoff.
These are statements with fixed point and \(j\), followed by constants,
followed by a single sufficiently small \(\eps_0\), and only then all
orders in the indicated window.



\section{The final selected contour}
\label{sec:F}

The selected upper index supplies a fixed product gap on a parameter
disk whose radius decreases only with particle number. Compact-root
pair contraction then controls the possible growth of the angular
Pfaffian. Cauchy estimates are applied to the complete weighted
integral on that disk.
\begin{proposition}\label{prop:F}
The bound \eqref{eq:reqF} holds for the final selected integral.
\end{proposition}
\begin{proof}
\noindent\textbf{The anchor and its parameter disk.} Telescope \(1-\psi\). A named selected index has \(p_q=1\), hence \(P\ge1\).
At \(\lambda=1\), this index is shifted inward by \(2\tau\).
The endpoint separation above gives
\(\Im W_q\ge\sinh(2\tau)\sin\delta\) at the radial center.
The interior root therefore stays in a strict subdisk, uniformly
on the selected support. It supplies a fixed \(Z\)-gap, and
\(\sum h_i\le-3\tau P/2\) supplies a fixed \(Y\)-gap.
Every true lower branch index is on \(m=1,p=0\) and has
\(\Im W\ge cP/N\), so its residue factor is at most \(C\sqrt N\).
Other roots are in a compact regular set or on the shifted upper
branch. The one-body product and pullback Jacobian cost
\(C^N N^{N/2}\).

These statements continue to the common disk \(|t-\se|\le c/N\):
the lower branch margin is \(c'/N\), compact roots move \(O(1/N)\),
and the fixed anchor absorbs the total compact-root inflation
\(\exp(CN(c/N))\) after the disk constant is made small.
The integral on this disk is holomorphic because the contour weights
are fixed in \(t\). Only the original undeformed representation needs
the smaller \(c\eps\) disk.

\par\smallskip\noindent\textbf{Root assignments and pair bounds.} The assignments in the following absolute estimates are made
pointwise at each fixed \(t\). Their domains need not be holomorphic
in \(t\), and no assigned integral will be differentiated.
Let \(M\) be the number of indices outside the true lower branch core.
Split their compact \(z\)-images into a strict interior group and a
compact lower-arc group. On this disk, same-group compact Schur
pairs have modulus at most \(q<1\); cross compact pairs are at most
\(1+C/N\); branch/compact pairs are also at most \(1+C/N\);
branch/branch pairs are at most one.

Here is the uniform mixed-pair calculation. On the entire protected
disk a true branch root \(w\) remains in the closed lower unit disk:
its \(W\)-image has strictly positive imaginary part by the
\(P/N\) margin. The compact radial-center root \(z_0\) belongs to
a fixed compact subset of that same half-disk excluding \(\pm1\).
Thus for some \(g>0\), either \(|z_0|\le1-g\) or
\(\Im z_0\le-g\). In both cases \(|1-z_0w|\ge g\).
For the second case this follows directly from
\[
 |z_0|^2|1-z_0w|^2
 =(|z_0|^2\Re w-\Re z_0)^2+
   (|z_0|^2\Im w+\Im z_0)^2.
\]
The compact root moves to \(z\) by at most \(L/N\); reduce the
disk constant so this is at most \(g/2\). With
\(A(z,w)=(z-w)/(1-zw)\),
\[
 A(z,w)-A(z_0,w)=\frac{(z-z_0)(1-w^2)}{(1-zw)(1-z_0w)},\qquad
 |A(z,w)|\le1+\frac{4L}{g^2N}.
\]
Only the compact root was displaced in this comparison: \(w\)
is held at its actual disk-point value. Hence no square-root branch
motion enters the estimate. Compact roots on both transition supports
are covered by the two compact groups, and no perturbed compact root
is assumed to lie in the unit disk. For compatibility with the existing
budget we may weaken this bound to \(1+C/\sqrt N\) in the next display.
Since the same-group count is at least \(M^2/4-M/2\),
\[
 |\Pf\widehat A(z)|\le
 \exp(CN+CM\sqrt N-aM^2)\le C_1^N e^{-a_1M^2}.
\]
The last inequality uses
\(CM\sqrt N\le(a/2)M^2+C^2N/(2a)\); the mixed inflation cannot be
dropped without this step.

Every \(y\)-pair with a true branch index is bounded. Indeed at
reciprocal angles its reflected upper partner has \(p=1\), and
the negative shift \(-2\tau\) dominates the possible positive
lower shift \(\tau P/(2N)\le\tau/2\).
For two compact angles write
\(d=\dist(\theta_i+\theta_j,2\pi\Z)\).
At exact reciprocity \(h_i+h_j\le0\) for every \(P/N\in[0,1]\);
smoothness gives \(h_i+h_j\le Ld\).
For \(d\le c\eps\) the negative base \(2\log r\) bounds the
modulus gap below by \(c'\eps\); for larger \(d\) the phase
bounds it below by \(c'd\). Thus
\[
 |A(y_i,y_j)|\le\frac C{\eps+d}.
\]
This majorant is independent of all other angles, including the
coupled scalar \(P\). Its matched two-angle integral is \(O(H+1)\).

\par\smallskip\noindent\textbf{Pfaffian summation and Cauchy estimates.} Write the augmented \(y\)-Pfaffian matrix as \(A_0+E\), putting the bounded ghost row in \(A_0\). The matrix \(E\) has entries
only between compact indices and \(A_0\) has bounded entries.
Expand in an even compact subset of size \(2r\). The remaining
Pfaffian minor is at most \(C^N N^{(N+1-2r)/4}\) by Hadamard.
There are at most \(N^{3r}\) coarse choices of subset and matching;
their disjoint kernels have integral at most \((C(H+1))^r\).
As \(r\le M/2\), a sufficient bound is
\[
 \int|\Pf\widehat A(y)|\le C^N N^{(N+1)/4}(CN^3(H+1))^M.
 \]
Combine these bounds and \(N!\ge(N/e)^N\) to obtain, for each assignment,
\[
 |F_{N,\mathrm{assignment}}(t)|
 \le C^N N^{-N/4+1/4}(CN^3(H+1))^M e^{-a_1M^2}.
 \]
At most $C^N$ assignments and polynomial named-index counts can be
absorbed in $C^N$. Complete the square in $M$ and use
\[
 [\log(CN^3(H+1))]^2\le C_2\bigl(1+\log^2(N+1)+\log^2(H+1)\bigr).
\]
The complete $F_N$ is holomorphic on a neighborhood of its closed
$c/N$ disk: the weights are fixed and the displayed integrable
majorant is uniform there. Cauchy's inequality, applied to the
complete integral and not to its pointwise assignments, gives
\[
 |F_N^{(j)}(t_\eps)|\le
 e^{B_j\log^2(H+1)}\,
 C_j^N N^{-N/4+j+1/4}e^{B_j\log^2(N+1)}.
\]
The last $N$-factor is eventually at most $N^{-N/8}$ and is
summable over all odd $N$, proving \eqref{eq:reqF}. All preceding
disk and pair bounds hold for every $N$: when $N\eps$ is large,
the branch one-body bound $C\sqrt N$ only becomes weaker, and
the compact-root inflation remains $e^{CN(c/N)}$. Thus no upper
particle-number window is used in this proposition.
\end{proof}



\section{The original contour and the all-branch region}
\label{sec:K}

This section treats the lower support left on the original contour.
The fixed partition identifies the branch, mixed, and compact regions.
Pair suppression and the original parameter disk give estimates for
large particle number. In the intermediate window, the all-branch
region is split into a nonresonant microcore and its complement; the
latter uses the fixed-domain calculus of Appendix~\ref{app:lie}.
\begin{lemma}[Fixed support partition]\label{lem:partition}
The original \(K\)-integral has a finite smooth partition into an
all-lower-branch region and regions of the following types: a compact
left root with a modulus gap, a separated compact/true-branch pair,
and an all-compact right region. The partition uses at most \(C^N\)
labels, and every finite cutoff jet retains its named anchor and its
stated support.
Each named small-\(\lambda\) Stokes current has the same three
compact/branch types, with no all-branch sector.
\end{lemma}
\begin{proof}
Use \(u=\theta+\theta_B\) near the lower branch. Fix a small
\(\delta_B>0\) and a large fixed \(L\).
Take \(\chi_B=1\) for \(|u|\le\delta_B/2\), supported in
\(|u|\le\delta_B\). Take \(\chi_{\mathrm{in}}=1\) for
\(|u|\le r_0=\delta_B/(2L^2)\), supported in \(|u|\le2r_0\).
The deformation function \(m\) is one on a neighborhood of these
supports.
Split \(1=\prod\chi_B+(1-\prod\chi_B)\) and telescope the second
term. It has a named \(q\) with \(|u_q|\ge\delta_B/2\),
including on the support of every derivative of that named factor.
Expand the other coordinates with \(1=\chi_{\mathrm{in}}+(1-\chi_{\mathrm{in}})\).

To make the left/right labels periodic, choose an open arc \(J\)
whose closure lies inside the upper \(a=1\) plateau. Cut the circle
at a point of \(J\) and unwrap the angle on its complement. Near
\(-\theta_B\), use \(u=\theta+\theta_B\) and a smooth right step
which is zero for \(u\le-r_0/4\) and one for \(u\ge r_0/4\).
Away from this local transition keep its value constant on each
side, and return smoothly from one to zero inside \(J\).
This defines a smooth periodic right label; the left label is its
complement.

At the branch transition every derivative of the products with
\(1-\chi_{\mathrm{in}}\) vanishes, since \(\chi_{\mathrm{in}}=1\) there. The additional
transition inside \(J\) is also harmless: \(1-a\) vanishes there
with all its jets. Thus the full \(K\)-weight and every current
weight \(\partial_{\theta_q}\psi\), together with all their jets,
vanish whenever any coordinate lies in \(J\). For the named current
coordinate this follows from \(a'=0\); for every other coordinate
it follows from its factor \(1-a\). On the support of any relevant
weighted jet the labels therefore select the required side.
On the transition of \(\chi_{\mathrm{in}}\) they are constant. Every effective
compact label excludes a fixed branch neighborhood.
If a true-branch label occurs, the outer
\(q\) and a true-branch \(j\) satisfy
\(|b_j|\ge4|b_q|\) after choosing \(L\) large, unless a left root
already supplies a modulus gap.
With no branch label, all roots are compact; either one is left or
all are right. Finite product differentiation creates at most
\(N^{O(j)}\) terms and preserves the supports just described.

For the current fix its named index \(q_S\) in
\(\operatorname{supp}a'\). Split its two upper endpoint arcs;
this split can be smooth where the current weight vanishes.
An upper-left \(q_S\) gives a fixed root modulus gap.
For an upper-right \(q_S\), expand every other index with
\(\chi_{\mathrm{in}}+(1-\chi_{\mathrm{in}})\) and the compact
left/right labels. A compact left index again gives a modulus
gap. Otherwise, if the true-branch label set \(J_B\) is nonempty,
use \(q_S\) and the fixed index \(\min J_B\) as the mixed pair;
if it is empty all indices are compact right. These alternatives
are exhaustive, including when all \(N-1\) other indices are
true branch. The named upper index prevents an all-branch sector.

The slope of \(q_S\) is uniformly bounded on a preliminary
endpoint neighborhood separated from the upper branch, independently
of subsequently shrinking the slivers. Choose the nested branch
margin so that its limiting slope dominates this bound by more
than four. Lemma~\ref{lem:branch}'s deformed estimates preserve
\(|b_{q_S}/b_j|\le1/4\) after the final choice of \(\eps_0\).
The named index remains on \(p=1,m=0\), and every true branch
on \(p=0,m=1\), on neighborhoods of all cutoff-jet supports.
The vanishing on \(J\) proved above applies also to the current.
Thus no cutoff derivative creates a new support type.
\end{proof}

\begin{lemma}[Pair suppression at radial centers]\label{lem:contraction}
After the fixed branch margins are chosen, the endpoint slivers and
deformation strength can be chosen so that at \(t=t_\eps\), on
\(K^N\) and on each Stokes support with one named \(q\),
\[
 \left|\prod_{i<j}P_{ij}\right|\le C^Ne^{-\kappa N^2},
 \qquad \kappa>0 .
\]
This is uniform for \(0\le\lambda\le1\), all permitted occupancies
\(P/N\in[0,1]\), and sufficiently small \(\eps\).
More precisely, lower indices can be assigned to three fixed
covering groups with same-group bound \(q'<1\) and cross-group
bound \(1+\eta\), where \(\log(1+\eta)<-\log(q')/4\).
Pairs containing the named Stokes index are bounded by \(C_S\ge1\).
For every set \(E\) of at most \(j\) pairs, these bounds imply
\[
 \prod_{\substack{i<l\\(i,l)\notin E}}|P_{il}|
       \le C_j^N e^{-\kappa N^2}.
\]
Fixed pair jets are bounded in regular branch coordinates and
compact angular coordinates; no uniform bound on raw branch-angle
jets is asserted.
\end{lemma}
\begin{proof}
\noindent\textbf{1. Fixed groups and limiting equality loci.}
Write \(b=\theta_B\), \(y_B=e^{-\ii b}\), and
\(\mu_0=S_*+1-\sqrt{(S_*+1)^2-1}\in(0,1)\).
Choose a fixed \(w<\min(b,\pi-b)/3\), small enough for the
regular lower-branch chart, with \([-b-w,-b+w]\) inside the
\(m=1\) neighborhood. Choose \(0<\ell_L,\ell_R<w\).
The closed limiting groups are
\[
 B=[-b-w,-b+w],\quad L=[-\pi,-b-\ell_L],\quad
 R=[-b+\ell_R,0].
\]
They cover the lower semicircle with overlaps. Assign each overlap
by any fixed rule. Add the upper-left sliver \([\pi-2\delta,\pi]\)
to \(L\) and the upper-right sliver \([0,2\delta]\) to \(R\),
identifying \(\pi\) and \(-\pi\). The named current index lies
in one of the two transition arcs. All cutoff jets retain these
closed supports. The branch group contains neither endpoint;
otherwise \(|P(y_B,1)|=|P(y_B,-1)|=1\) would prevent strictness.

At the limit \(W_0(\theta)=S_*-\cos\theta\) and
\(W_0+1\ge S_*>0\). On \([-\pi,-b]\) the roots are real
in \([\mu_0,1]\); on \([-b,0]\) they are \(e^{-\ii v}\),
\(0\le v\le b\). The identity
\[
 |1-zw|^2-|z-w|^2
 =(1-|z|^2)(1-|w|^2)+4\Im z\Im w
\]
and its unit-circle sine version give \(|P|\le1\).
Apart from the removable configurations described below, equality
requires an endpoint in the \(y\)-pair and equality in the
\(z\)-factor. Its complete loci, with either ordering, are
\(y_i=1,\ \theta_j\in[-\pi,-b]\), and
\(y_i=-1,\ \theta_j\in[-b,0]\). Common-endpoint pairs have
canceled value zero. No equality locus lies in a same-group square.
On \(L^2\) the roots stay in a strict real subinterval of \((0,1)\);
on \(R^2\) they stay on a compact lower arc away from the real axis.
On \(B^2\) the \(y\)-factor is bounded by
\[
 \frac{\sin w}{\min\{\sin(b-w),\sin(b+w)\}}
 \le\frac1{2\cos w}<1.
\]
The canceled extension below covers its double branch. Thus
the three compact same-group suprema have a maximum \(q<1\).
Cross-group suprema may equal one.

\par\smallskip\noindent\textbf{2. Local continuation and the complete pair.}
Use thin simply connected neighborhoods of the compact real
\(W_0>1\) intervals and of the compact intervals inside \((-1,1)\),
avoiding both branch points. On the former choose the root real
in \((0,1)\); on the latter choose the continuation of
\(e^{-\ii\arccos W_0}\). These local roots may leave the unit disk.
Near the lower branch use \(\phi=\arccos W\) on the
fourth-quadrant sheet for \(\Im W\ge0\).
The implicit-function theorem gives the even analytic inverse
\(y_t(\phi)=G_t(\phi^2)\) near \((t_*,0)\), since
\(\partial_yW(y_B;t_*)\ne0\). In this neighborhood
\[
 y_t(\phi_i)-y_t(\phi_l)
  =(\phi_i-\phi_l)(\phi_i+\phi_l)H_t(\phi_i,\phi_l),\qquad
 1-e^{-\ii(\phi_i+\phi_l)}=(\phi_i+\phi_l)J(\phi_i+\phi_l),
\]
where \(J(0)=\ii\). Substitution into \eqref{eq:pair} gives
\(P_{il}=(\phi_i-\phi_l)^2A_t\), with analytic bounded fixed
jets and \(A_{t_*}(0,0)=-1/(4\sin^2b)\).
This resolves the whole local reciprocal-root hypersurface,
without dividing by a collision zero.

At the limiting lower support \(1-z_i z_l=0\) forces both
roots to be \(1\) or both \(-1\). The latter is excluded by
\(W_0+1\ge S_*\); the former is the resolved lower double branch.
On the closed compact complement of a smaller double-branch
neighborhood, the denominator therefore has a positive minimum
\(\mu\). Shrink its neighborhoods to retain \(\mu/2\).
The regular and rational expressions agree on overlaps by
\eqref{eq:pair}. This finite cover supplies continuous complete
pairs and bounded fixed jets in regular branch and compact
coordinates. No globally glued square-root domain is needed.
The apparent \(1-y_i y_l\) pole was already canceled; the upper
inverse branch, which would make a genuine pole with the lower
branch, is excluded by the selector plateau.

For a named right limiting endpoint root \(\xi=e^{-\ii b}\),
\(|1-\xi z|=|e^{\ii b}-z|\ge\sin b\) for every limiting lower
root. For the left endpoint it is at least \(1-\mu_0\).
Thus a positive named margin
\(\mu_S=\min(\sin b,1-\mu_0)\) persists, with \(\mu_S/2\)
on sufficiently small neighborhoods. Boundedness of the canceled
numerator supplies \(C_S\ge1\); no named contraction is needed.

\par\smallskip\noindent\textbf{3. Uniform support transfer and reserved slack.}
Project each sliver to its endpoint and keep each lower angle
fixed; denote the projected angle by \(\bar\theta_i\).
For \(\rho=P/N\in[0,1]\) and \(0\le\lambda\le1\),
\[
 v_i=\log|y_i|=-c_0\eps+\lambda\tau(-2p_i+\rho m_i/2),
 \qquad |v_i|\le c_0\eps+5\tau/2.
\]
Let \(L_S\) bound \(|S'|\) on a fixed parameter neighborhood.
Uniformly in this entire parameter box, hence for every actual
coupled occupancy and every \(N\),
\[
 |y_i-e^{\ii\bar\theta_i}|+|W_i-W_0(\bar\theta_i)|
 \le C(\eps+\tau+\delta)+L_S|t-t_\eps|.
\]
Compact-root changes are Lipschitz. Near the branch, the factorization
\(\phi=\sqrt{2(1-W)}h(1-W)\), with analytic \(h\), gives a
one-half Holder bound on the chosen closed half-plane.
The regular pair is Lipschitz in \((t,\phi_i,\phi_l)\);
on the compact complement the denominator margin gives the same
property in \((y_i,y_l,z_i,z_l)\). A finite cover therefore gives
one uniform complete-pair modulus, bounded by
\(C\sqrt{\eps+\tau+\delta+|t-t_\eps|}\) in these neighborhoods.
It applies whenever the branch-near roots stay on the designated
sheet; the two disk proofs below verify this condition separately.

Choose \(q<q_1<q'<1\), then \(\eta>0\) with
\(\log(1+\eta)<-\log(q')/4\).
Choose the slivers and \(\tau\), and a preliminary \(\eps_0\),
so the radial-center modulus is smaller than both \(q_1-q\)
and \(\eta/2\), while retaining the denominator neighborhoods.
This gives same-group bound \(q_1\) and cross-group bound
\(1+\eta/2\); reserve \(q'-q_1\) and \(\eta/2\) for the disks.
The real radial motions increase \(\Im W\), so every radial-center
root is interior. This transfer includes the slivers in the
same-group estimate as well as the cross estimate.

For clarity, the limiting branch/upper-sliver moduli for small
\(a>0\) are
\[
 \frac{\sin((b+a)/2)}{\sin((b-a)/2)}
   =1+a\cot(b/2)+O(a^2),\qquad
 \frac{\cos((b-a)/2)}{\cos((b+a)/2)}
   =1+a\tan(b/2)+O(a^2).
\]
Holding a sliver fixed while moving a compact-group boundary
toward the branch can consequently make an \(RR\) or \(LL\)
pair exceed one. The margins must therefore be fixed first, the contraction slack second,
and the sliver width last.

Let \(n\) be the number of lower variables: \(n=N\) for \(K\)
and \(n=N-1\) for a Stokes term. Assign each to one covering group.
After at most \(j\) pairs have been removed, the remaining
same-group count is at least \(n^2/6-n/2-j\), while the
remaining cross-group count is at most \(n^2/3\). Apply the
individual upper bounds directly to the remaining product:
\[
 (q')^{n^2/6-n/2-j}(1+\eta)^{n^2/3}
       \le(q')^{n^2/12-n/2-j}.
\]
This remains valid if the lower bound on the count is negative.
It involves no division by a removed factor, which may vanish.
The named upper compact
\(q\) has at most \(N-1\) additional pairs, each uniformly
bounded: its \(z_q\) is separated from the inverse of a lower
root, including the branch root. These pairs cost \(C_S^{N-1}\).
For \(n=N-1\) the displayed exponent is
\(N^2/12-2N/3+7/12-j\); the linear correction is included in
\(C_j^N\). Thus \(\kappa=-\log(q')/12\) works for both objects.
The canceled double-branch expression is analytic in the regular
coordinates, and the other pair denominators are separated on the
fixed compact sets. This also proves the stated fixed-jet bounds.
\end{proof}

\begin{lemma}[The original parameter disk]\label{lem:originaldisk}
For a sufficiently small fixed \(c>0\), the original contour
density and all \(K\) and Stokes transition densities are holomorphic
on \(|t-t_\eps|\le c\eps\). On the \(K\) and named Stokes
supports, the following estimates hold:
\[
 |1-\sigma Y|,\ |1-\sigma Z|\ge c'\eps,\qquad
 \left|\prod_{i<j}P_{ij}\right|\le C^Ne^{-\kappa' N^2}.
\]
The integrated one-body majorants and contour Jacobian cost
at most \(C^N\), uniformly in \(N\) and \(0\le\lambda\le1\).
\end{lemma}
\begin{proof}
\noindent\textbf{Product gaps.} The undeformed dispersion has \(\Im W\ge c_2\eps\).
Every permitted radial motion increases this quantity, as in
Proposition~\ref{prop:selector}; reducing \(c\) preserves a
fixed fraction of the margin on the disk. Each root is then
in the lower unit half-disk with attenuation at least \(c_3\eps\).
Thus the product gaps are bounded below by a constant times
\(\min(1,N\eps)\), which is at least a constant times \(\eps\).
The contour \(Y\) is independent of \(t\).

\par\smallskip\noindent\textbf{An integrable branch majorant.} Near the lower branch put \(\mu=\lambda P/N\),
\(E=\eps+\mu\), and \(\delta=|t-\se|\le c\eps\).
On its true plateau, write
\(y=e^{v+\ii\vartheta}\), where
\(v=-c_0\eps+\tau \mu/2\) and \(\vartheta=-\theta_B+u\).
The exact identities
\[
 \Re W=\Re S-\cosh v\cos\vartheta,\qquad
 \Im W=\Im S-\sinh v\sin\vartheta
\]
give, with \(D_B=1-W\) and \(a_B=\sin\theta_B>0\),
\[
 |\Re D_B-a_Bu|
 \le C(u^2+\eps^2+\mu^2+\delta),\qquad
 \Im W\ge c_bE-C_S\delta\ge(c_b/2)E.
\]
Here the last inequality holds after reducing the disk constant
\(c\). All constants depend only on the fixed point and margins.
In particular, the real branch-center displacement on this disk is
\(O(\eps^2+\mu^2+\delta)\): the parameter perturbation cannot be
omitted from the radial-center estimate.

Choose a fixed branch radius \(h\) and a fixed \(K\ge1\).
If \(|u|\le KE\), the imaginary margin gives
\(|D_B|\ge c_b(|u|+\eps)/(2(K+1))\).
If \(KE<|u|\le h\), then
\(u^2\le h|u|\), \(\eps^2+\mu^2\le E^2\le h|u|/K^2\),
and \(\delta\le cE\le c|u|/K\).
Choose \(h,K,c\) so that
\(C(h+h/K^2+c/K)\le a_B/2\), in addition to the previous margins.
Then \(|\Re D_B|\ge a_B|u|/2\) in this second region.
These choices are independent of \(N,\eps,\lambda,P\).
Since \(W+1\) stays separated from zero on the fixed branch arc
and \(R=z/\sqrt{W^2-1}\), both regions yield
\[
 |R(y;t)|\le \frac{C}{\sqrt{|u|+\eps}}
\]
uniformly on the full disk. The upper branch is absent from
these supports; the remaining compact factors are bounded.
The displayed majorant does not contain the coupled quantity
\(P\), and its integral on \(|u|\le h\) is at most \(4C\sqrt h\).
Thus its product can be integrated with cost \(C^N\) without
treating \(P\) as constant. The inverse-product numerator
and the Jacobian in Proposition~\ref{prop:selector} also cost \(C^N\).

\par\smallskip\noindent\textbf{Persistence of pair contraction.} For an explicit disk choice, let \(b_0\eps\) be the
undeformed imaginary margin and let \(L_S\) be as in
Lemma~\ref{lem:contraction}. Choose \(L_Sc\le b_0/4\).
The monotone radial motions then leave \(\Im W\ge3b_0\eps/4\)
throughout the disk. Along each disk radius,
\(\partial_W\phi=-1/\sin\phi\) gives branch-coordinate motion
\(O(c\sqrt\eps)\); compact roots move by \(O(c\eps)\).
The finite continuation cover and the regular canceled expression
of Lemma~\ref{lem:contraction} give a uniform complete-pair change
\(O(\sqrt\eps)\). Reduce the final \(\eps_0\) to make it
smaller than both reserved margins \(q'-q_1\) and \(\eta/2\).
Thus same pairs are at most \(q'\), cross pairs at most
\(1+\eta\), and named pairs are uniformly bounded.
The limiting cross supremum need not be strict. The direct
remaining-pair count proves the estimate with the same
\(\kappa=-\log(q')/12\), or a smaller common exponent.
The pair estimates apply to the stated \(K\) and current supports,
where the upper inverse branch is excluded. This support restriction
is what supplies the pair contraction. The original-disk proof is
independent of the large-\(\lambda\) estimate.
\end{proof}

\begin{lemma}[Branch geometry]\label{lem:branch}
On a sufficiently small original lower-branch arc,
\[
 1-W=a_Bu-\ii b_B\eps+O(u^2+\eps|u|+\eps^2),
 \quad a_B=\sin\theta_B>0,\quad b_B=2\sin\theta-c_0a_B>0.
\]
With \(\phi=\arccos W\) on the fourth-quadrant sheet,
\begin{align*}
 |\phi'(u)|&\asymp(|u|+\eps)^{-1/2},&
 \int_B|\phi'(u)|\dd u&\le C,\\
 -\Im\phi(u)&\ge c\sqrt{|u|+\eps}\quad(u\le0),&
 x'(u)&\asymp(u+\eps)^{-1/2}\quad(u\ge0),
\end{align*}
where \(x=\Re\phi\).
For \(u\ge C_0\eps\), \(-x''(u)\ge c u^{-3/2}\).
If \(M/\eps\to\infty\) uniformly and \(0\le m\le M/2\),
\[
 x'(m)-x'(M)\ge c(m+\eps)^{-1/2},\qquad
 \frac{|\phi'(m)\phi'(M)|}{x'(m)-x'(M)}\le C M^{-1/2}.
\]
If \(M/2<m<M\) and \(m\gg\eps\), the latter ratio is at most
\(C\sqrt M/(M-m)\).

On a true-branch plateau of the current put
\(\mu=\lambda P/N\), \(E=\eps+\mu\). Uniformly for
\(0\le\lambda\le\lambda_*\) and \(|u|\le2r_0\),
\begin{align*}
 |\phi_u|&\asymp(|u|+E)^{-1/2},\\
 -\Im\phi&\ge c\sqrt{|u|+E}\quad(u\le0),&
 \Re\phi_u&\ge\tfrac23|\phi_u|\quad(u\ge0).
\end{align*}
The constants may depend on the fixed \(\tau\), but not on
\(N,\eps,\lambda,P/N\). Here the margins and then \(\eps_0\)
are small enough, and \(P\) is fixed during a plateau variation.
\end{lemma}
\begin{proof}
The displayed normal form follows by Taylor expansion of \(S\) and
the fixed-radius \(y\). Put \(D=1-W\).
The exact factorization \(\phi=\sqrt{2D}\,h(D)\), \(h(0)=1\),
with analytic \(h\), yields the first estimates and an integrable
branch length. On \(u\ge C_0\eps\), the leading terms
\(x'(u)=\sqrt{a_B/2}\,u^{-1/2}(1+O(\eps/u+u))\) and
\(-x''(u)\asymp u^{-3/2}\) have the stated signs after increasing
\(C_0\) and shortening the arc.
If \(m\ge C_0\eps\) and \(m\le M/2\), the relative error in
\(x'\) can be made smaller than a fixed fraction of
\(1-1/\sqrt2\); this gives the strong difference estimate.
If \(m<C_0\eps\), its slope is at least \(c\eps^{-1/2}\),
whereas the slope at \(M\) is \(o(\eps^{-1/2})\).
For comparable \(m,M\), integrate the bound on \(-x''\).
The sector containing \(\phi'\) also has width less than \(\pi\);
integrating in \(u\) gives
\(|\phi(u)-\phi(v)|\ge c|u-v|\).

For the current, write \(y=e^{v+\ii\vartheta}\),
\(v=-c_0\eps+\tau \mu/2\), \(\vartheta=-\theta_B+u\).
Put \(Q=\eps+\tau \mu\), comparable to \(E\) for fixed \(\tau>0\).
The exact real and imaginary identities in
Lemma~\ref{lem:originaldisk}, at the radial center, give
\begin{gather*}
 D=1-W=a_Bu-\ii\beta+R,\qquad cQ\le\beta\le CQ,\\
 R\in\R,\qquad |R|\le C(u^2+Q^2),\qquad
 D_u=a_B+O(|u|+Q).
\end{gather*}
Positivity of \(\beta\) uses \(b_B>0\) and the positive
\(\tau \mu\) coefficient. The constants in these \(Q\)-bounds
can be chosen before shrinking \(\tau\). The model
\(a_Bu-\ii\beta\) has modulus comparable to \(|u|+Q\), and
its relative perturbation is \(O(r_0+Q)\). Its square root for \(u\le0\) has argument
in \([-\pi/2,-\pi/4]\); its inverse square root for \(u\ge0\)
has argument in \([0,\pi/4]\). Using the exact factorization
\(\phi=\sqrt{2D}\,h(D)\) and its derivative proves the asserted
cone and slope bounds after shrinking the fixed margins and
\(\eps_0\). This argument allows the \(O(E^2)\) real displacement
of the branch center; it does not require \(\Re D\le0\) at \(u=0\).
In particular, since all \(|z_i|\le1\), for a negative branch
coordinate
\[
 \frac{|\phi'_v|}{|1-\sigma Z|}
 \le\frac{|\phi'_v|}{1-|Z|}
 \le\frac C{|u_v|+\eps+\lambda P/N}.
\]
\end{proof}

\begin{proposition}[Nonresonant microcore]\label{prop:micro}
For \(N_0+2\le N<D\sqrt H\), put
\[
 b_N=c_*e^{-(A+4)N}/N,\qquad
 \psi_N=\prod_i\chi(u_i/b_N),
\]
where \(\chi=1\) on \([-1/2,1/2]\), supported on \([-1,1]\), and
\(c_*>0\) is fixed and small. The all-\(B\) microcore contribution
obeys, for every \(j\le k\),
\[
 |\partial_t^j I_{\mathrm{micro},N}|
 \le \exp(C_jN^2)\,b_N^{(N^2-1)/2}.
\]
\end{proposition}
\begin{proof}
First use the intermediate particle window to compare the radial and
microcore scales. Uniformly in this window,
\(N\eps e^{AN}\to0\) and \(\eps/b_N^m\to0\) for every fixed \(m\).
Thus on the microcore and all its cutoff jets,
\[
 |Y-\xi^N|\le CN(\eps+b_N)\le\tfrac12e^{-AN},
 \qquad |1-\sigma Y|\ge\tfrac12e^{-AN}.
\]
The one-phase field \(V_i=a_i=S'/W_{u_i}\) freezes each regular
\(\phi_i\), hence \(Z\). In these coordinates the numerator is
analytic times \(\Delta(\phi)^2\).
Repeated Lie differentiation leaves the latter polynomial unchanged
and leaves the \(Z\)-kernel simple. The \(Y\)-power is at most
\(j+1\); its gap and the \(b_N^{-j}\) cutoff cost are
\(\exp(O_j(N))\). Analytic products have bound \(\exp(O_j(N^2))\).

The branch image lies in a fixed cone of width less than \(\pi\)
and has radius \(C\sqrt{b_N}\). Decreasing \(c_*\) ensures
\(N\sqrt{b_N}\) is small, so $|1-Z|\ge c\sum_i|\phi_i|$. If $\sigma=-1$,
$Z$ is close to one, and $|1+Z|\ge1$; since $\sum_i|\phi_i|$
is small, the same weaker lower bound works for either sign.
The square-root normal form gives, directly on both sides of zero,
\(|\phi|\ge c\sqrt{|u|+\eps}\) and
\(|\phi'|\le C(|u|+\eps)^{-1/2}\). Substituting
\(w=\sqrt{|u|+\eps}\) therefore gives
\[
 \int_{\mathrm{core}}e^{-\mu|\phi|}|\dd\phi|
       \le4C\int_{\sqrt\eps}^{\sqrt{b_N+\eps}}e^{-c\mu w}\dd w
       \le C\min(\sqrt{b_N},\mu^{-1}).
\]
Using \(1/\sum|\phi_i|=\int_0^\infty e^{-\mu\sum|\phi_i|}\dd \mu\)
therefore costs \(C^N b_N^{(N-1)/2}\).
The Vandermonde costs at most
\(C^{N(N-1)}b_N^{N(N-1)/2}\). These estimates prove the claim.
The logarithm of the displayed majorant is
\[
 -\frac{A+4}{2}N^3-\frac12N^2\log N+O_j(N^2).
\]
In particular it is at most \(-(A+4)N^3/2+O_j(N^2)\).
Thus the sum of these bounds over all \(N\ge N_0+2\) is finite.
\end{proof}

\begin{proposition}[The all-branch complement]\label{prop:allB}
The remaining original all-\(B\) contribution, for
\(N_0+2\le N<D\sqrt H\), is bounded through every \(j\le k\) by
\[
 \exp(C_jN\log(N+1))(H+N)^{C_j}.
\]
\end{proposition}
\begin{proof}
\noindent\textbf{1. Fix the anchor and the shape scale.} Telescope \(1-\psi_N\); each term has a named anchor \(a\) with
\(|u_a|\ge b_N/2\), including all cutoff jets.
Let \(\mu=N^{-1}\sum u_i\), \(\rho^2=\sum(u_i-\mu)^2\), and split
smoothly at \(\rho_N=e^{-BN}\), using a function of
\(\rho^2/\rho_N^2\). Put
\[
 \Gamma_* = \sup_{\substack{n\in\mathbb N\\n\ge N_0+2}}
                   \frac{\log(16n/c_*)}{n}<\infty.
\]
Choose \(B>A+4+\max\{1,\Gamma_*\}\), increasing it as needed
after the finite-jet constants have been fixed. Then
\[
 2\rho_N\le b_N/8\qquad\hbox{for every }N\ge N_0+2.
\]
This single choice gives the support inclusion at every order,
including the finitely many smallest orders.

\par\smallskip\noindent\textbf{2. Integrate the far region.} On the far support \(\rho\ge\rho_N\), the maximal separation
is at least \(\rho_N/\sqrt N\). The rational two-phase Lie
calculus costs \(\exp(O_j(N\log N))\), including every exponential
cutoff jet. The undifferentiated pairs retain
\(\exp(-\kappa N^2+O_j(N))\).
It remains to integrate two simple global kernels against
\(\prod_i|\phi'_i|\dd u_i\).

If all coordinates are nonnegative, their maximum \(M\ge b_N/2\).
Select the actual minimum \(m\) and maximum \(M\) only now,
after differentiation. For fixed other variables the real phase
map is
\[
 (m,M)\longmapsto(F,G)
    =(m+M+\mathrm{const},\,x(m)+x(M)+\mathrm{const}).
\]
Its Jacobian \(x'(m)-x'(M)>0\) makes it injective on each fixed-\(F\)
extreme-index cell. Lemma~\ref{lem:branch} bounds the ratio of its
two branch measures to that Jacobian by \(\exp(O(N))\):
use \(M\ge b_N/2\) for \(m\le M/2\), and
\(M-m\ge\rho_N/\sqrt N\) otherwise.
Only \(m,M\) vary here; all other angles are fixed. Each unwrapped
phase-image interval has bounded length, independent of \(N\), since
\(m,M\) and \(x(m),x(M)\) range over fixed bounded intervals.
The additive constants translate these intervals without changing
their lengths. Thus each simple periodic phase integral costs
\(O(H+N)\). The other branch lengths cost \(C^N\).

For negative coordinates choose \(\delta_N=e^{-C_1N}\) with
\(C_1>A+4+\max\{1,\Gamma_*\}\). Then \(\delta_N\le b_N/16\)
for every \(N\ge N_0+2\). If some negative \(|u_v|\ge\delta_N\), then
\(|1-\sigma Z|\ge e^{-O(N)}\); integrate the remaining \(Y\)-kernel
in the named anchor, whose measure is at most \(e^{O(N)}\dd u_a\).
If all negative magnitudes are smaller than \(\delta_N\), the
named anchor is positive. Choose one negative \(v\); the bound
\[
 \frac{|\phi'(u_v)|}{|1-\sigma Z|}\le\frac C{|u_v|+\eps}
\]
has integral \(O(H+N)\). The positive anchor integrates the
\(Y\)-kernel with cost \(e^{O(N)}(H+N)\).
The scale inclusion already covers every finite order; any
remaining finite-order constants only enlarge the numerical bound.
This proves the far bound
\(\exp(-\kappa N^2+O_j(N\log N))(H+N)^2\).

\par\smallskip\noindent\textbf{3. Use the collision zero in the near region.} On the near support, \(|u_i-\mu|\le\rho\le2\rho_N\le b_N/8\).
The anchor \(|u_a|\ge b_N/2\) gives \(|\mu|\ge3b_N/8\), so
all coordinates have the same sign and \(|u_i|\ge b_N/4\).
The pointwise branch-measure product has the bound
\[
 \prod_i|\phi'_i|\le(Cb_N^{-1/2})^N
 =\exp\left(\frac{A+4}{2}N^2+
             \frac N2\log(N/c_*)+O(N)\right).
\]
Segments joining two coordinates stay on this same-sign support,
so each collision difference is at most \(Cb_N^{-1/2}\rho\).
The collision coefficient and the branch measures together are
therefore bounded by \(C^{N^2}b_N^{-N^2/2}\), whose logarithm is
\[
 \frac{A+4}{2}N^3+\frac{N^2}{2}\log(N/c_*)+O(N^2).
\]
Fixed-\(j\) coefficient and derivative costs can be included in
the lower-order \(O_j(N^2)\) term. Its constant is fixed before
\(B\). These bounds are determined
on the \(B\)-independent anchor neighborhood; the positive cubic
coefficient is independent of \(B\). Lie and cutoff jets lose at most
\(2j+1\) powers of \(\rho\).
The shape measure contributes \(\rho^{N-2}\).
For negative \(\mu\), the \(Z\)-gap is \(e^{-O(N)}\) and one mean
phase integral suffices.
For positive \(\mu\), strict concavity and zero first shape derivative
give \(-\partial_\rho\Re\sum\phi_i\ge c\rho\), while the \(Y\)-phase
is \(N\mu+\mathrm{const}\). Their real Jacobian is at least \(cN\rho\).
After coarea the conservative remaining power is
\[
 L=N^2-2j-4\ge N^2-2k-4\ge4N_0+1>0.
\]
Bound this power by \((2\rho_N)^L\) before integrating the phase
kernels. At fixed shape direction, \(F=N\mu+\mathrm{const}\) fixes
\(\mu\), and strict monotonicity of \(G\) in \(\rho>0\) gives
multiplicity one for the unwrapped map. Each phase lies in an
interval of length at most \(CN\): the mean angle ranges over a
fixed interval and each individual branch phase is bounded.
The periodic-kernel estimate in Appendix~\ref{app:lie} therefore
costs at most \(CN^2(H+N)^2\) for the two phases, rather than a
single-period bound. For negative \(\mu\), the sole mean phase
instead costs at most \(CN(H+N)\), with the separated \(Z\)-kernel
already in the preceding estimate. These polynomial factors are
absorbed in \(e^{C_jN\log(N+1)}\). At this stage the shape radius has been exchanged for a phase
variable. With its remaining power already bounded on the support,
the bound is at most
\[
 \exp\left(\frac{A+4}{2}N^3+\frac{N^2}{2}\log N+C_jN^2\right)
 (2\rho_N)^L(H+N)^2.
\]
Its logarithm, apart from the simple-kernel factor, is at most
\(((A+4)/2-B)N^3+\tfrac12N^2\log N+O_j(N^2)+B(2j+4)N\).
The already required \(B>A+4\) makes this tend to minus infinity.
Finitely many remaining \(N\)'s cost a fixed, possibly very large,
\(\eps\)-independent constant.
At fixed \(\eps>0\) the same positive collision power makes the
full-equality puncture flux vanish at every preceding Lie stage.
Thus the identities used on the punctured charts have no residual
boundary term. The near and far bounds prove the proposition.
\end{proof}



\section{Compact and mixed sectors; Stokes estimates}
\label{sec:sectors}

The remaining intermediate-order sectors have either a compact root
with a modulus gap or a pair of coordinates that controls the two
product phases. These give the first two propositions. The last two
estimates use holomorphic parameter disks for the large-current
interval and the high-particle window, respectively.
\begin{proposition}[Separated mixed and left sectors]\label{prop:mixed}
For \(N_0+2\le N<D\sqrt H\), the non-all-\(B\) original \(K\)
sectors with a compact left root or a separated compact/branch
pair, and the corresponding small-\(\lambda\) Stokes sectors,
obey \eqref{eq:reqinter}. The differentiated current density
obeys the per-order bound uniformly for \(0\le\lambda\le\lambda_*\).
\end{proposition}
\begin{proof}
\noindent\textbf{A field that freezes the product phases.} Use hybrid coordinates consisting of regular \(\phi_i\) for true
branch indices \(J\) and original angles for the compact indices.
All coefficient jets are bounded in these coordinates.
For Stokes the named \(q\) lies on \(p=1,m=0\), and every true
branch lies on \(m=1,p=0\); their angle variations preserve both
\(P=\sum p_i\) and \(\sum m_i\) exactly.
For original \(K\) the deformation is zero, so these occupancy
constraints are unnecessary.

Put \(a_i=(\partial_t\phi_i)/b_i\),
\(b_i=\partial_{\theta_i}\phi_i\),
\(A=\sum_{i\in J}a_i\), and
\(D=\sum_{i\notin J}\partial_t\phi_i\).
With \(j=\min J\) fixed for the labelled term, let
\[
 X=\frac{D+b_qA}{b_j-b_q},\qquad
 V_i=a_i\ (i\in J\setminus\{j\}),\quad
 V_j=a_j+X,\quad V_q=-A-X .
 \]
All other velocities vanish. Direct substitution freezes \(Y,Z\)
and every unselected branch phase. The residual branch velocity is
\[
 (\partial_t-V\cdot\nabla)\phi_j
       =-\frac{D+b_qA}{1-b_q/b_j}.
\]
The support choice gives \(|b_q/b_j|\le1/4\).
Since \(1/b_j\) is analytic and vanishes in its regular branch
coordinate, this residual, its finite jets, and the compact
velocity have only polynomial \(N\)-cost, uniformly in \(\eps,\lambda\).
Here the hybrid top form can be written explicitly. Up to the
unchanged factor \((Z^{-1}+Y^{-1})/[N!(1-\sigma Z)(1-\sigma Y)]\prod_{i<l}P_{il}\),
its measure is
\[
 \prod_{i\in J}\left[-\frac{z_i}{\pi(1-y_i^{-2})}\dd\phi_i\right]
 \prod_{i\notin J}\left[\frac{R_i y_i d_i}{2\pi\ii}\dd\theta_i\right].
\]
Indeed \(d_i=\ii\) on every true branch plateau, the angular
determinant is \(\prod d_i\), and the hybrid change of variables
has triangular determinant \(\prod_{i\in J}b_i\).
This retains the full Jacobian of the coupled contour.

For Stokes, \(V\cdot\nabla P=V\cdot\nabla\sum m_i=0\) identically
on an open neighborhood of the support. The same is true for
all parameter derivatives of this identity. Thus the coupled
occupancy introduces no hidden branch motion at higher order.
The actual current multiplier is
\(-2\ii\tau\,\partial_q\psi\); its finite derivatives, and the
derivatives of the nested partition, stay on these plateaus.
Use the pullback recurrence \eqref{eq:pullbackLie} in the hybrid
coordinates. Its residual coefficients have no difference pole,
because \(1-b_q/b_j\) is separated from zero. Each fixed operation
changes finitely many analytic pair or one-body factors and
introduces at most polynomial \(N\)-cost. The remaining pair
factors therefore retain
\(C_j^N N^{C_j}e^{-\kappa N^2}\).

\par\smallskip\noindent\textbf{Integrating the remaining simple kernels.} For positive branch coordinates, Lemma~\ref{lem:branch}, including
its current version, gives \(\Re b_j\ge(2/3)|b_j|\).
The compact slope is dominated by the true branch slope; for a
Stokes upper-right \(q\) it has the opposite sign in the limit.
The phase map
\((\theta_q,u_j)\mapsto(\sum\theta_i,\Re\sum\phi_i)\)
has Jacobian \(\Re(b_j-b_q)\ge(2/3-1/4)|b_j|\)
and is injective on fixed first-phase slices.
It absorbs the selected branch measure, and each remaining phase
kernel integrates with cost \(O(H+\log N)\).
Every unselected branch measure has bounded integral.

If a branch coordinate \(u_v<0\), its attenuation and branch measure
give instead
\[
 \frac{|\phi'_v|}{|1-\sigma Z|}
 \le\frac C{|u_v|+\eps+\lambda P/N}.
\]
Integrate this factor in \(u_v\), and the \(Y\)-phase in the compact
\(q\). Both cost at most \(C(H+\log N)\).
The phase statements use variables on true plateaus; therefore \(P\)
does not vary during the selected integrations.
These sign divisions are performed only on the final absolute integral.

If a compact left root occurs, \(Z\) has a fixed modulus gap.
Set \(V_i=a_i\) for branches and \(V_q=-A\), preserving \(Y\)
and freezing every branch root. Differentiating the separated
\(Z\)-kernel costs only polynomial \(N\), and \(q\) integrates
the simple \(Y\)-kernel. A left \(q\) itself is allowed.
The fixed support partition exhausts the asserted cases.
\end{proof}

\begin{proposition}[Compact right sectors]\label{prop:compactR}
The all-compact right original \(K\) and small-\(\lambda\) Stokes
sectors satisfy the same per-order bound in \eqref{eq:reqinter}.
\end{proposition}
\begin{proof}
\noindent\textbf{Curvature and the selected field.} On a fixed compact right interval \(R\), separated from both branch
points, put \(g(\theta)=\arccos(S_*-\cos\theta)\). Direct differentiation
gives
\[
 g''(\theta)=-\frac{S_*(1-\cos\theta\cos g(\theta))}{\sin^3g(\theta)}
            \le-c_R<0 .
\]
The coupled contour is smooth in the angles and analytic in \(t\).
Its fixed finite angular jets differ from their limiting values by at most
\(C_jN^{C_j}(\eps+\lambda)\). Indeed \(P/N\in[0,1]\), each
fixed angular jet of \(P/N\) is uniformly bounded, and every
deformation term carries a factor \(\lambda\). The undeformed
parameter differs from \(t_*\) by \(O(\eps)\). The root map has
bounded jets on the fixed regular compact neighborhood; the chain
and product rules and the sum over indices cost only a polynomial
in \(N\) for fixed \(j\). For the Hessian, passing from entrywise
bounds to operator norm costs at most one further factor \(N\).
Since \(N\le D\sqrt H\), \(\eps=e^{-H}\), and
\(0\le\lambda\le e^{-\alpha_0H}\), these errors tend to zero
uniformly after all jet constants and \(D,\alpha_0>0\) are fixed.
In particular
\(\nabla^2 G\le-(c_R/2)I\) for \(G=\Re\sum\phi_i\).

Write \(f_i=\partial_{\theta_i}\log Y\), \(g_i=\partial_{\theta_i}\sum\phi_i\),
and \(Q_t=\partial_t\sum\phi_i\). Permutation symmetry makes
\(D_{pq}=f_pg_q-f_qg_p\) zero at \(\theta_p=\theta_q\).
Division by the difference is justified for these smooth angular
functions by the integral of their transverse derivative; no analytic
extension of the deformation bumps is assumed.
At the limit it equals \(\ii[g'(\theta_q)-g'(\theta_p)]\).
The strict curvature and the finite-jet perturbation bound yield
\[
 D_{pq}=(\theta_p-\theta_q)C_{pq},\qquad |C_{pq}|\ge c_R/2 .
\]
The field
\(V_p=-Q_tf_q/D_{pq},\ V_q=Q_tf_p/D_{pq}\)
freezes \(Y,Z\). Its simple difference pole, divergence, and
rational real-angle weights have the two-power-per-Lie-order
bound of Appendix~\ref{app:lie}.

\par\smallskip\noindent\textbf{Separated and near-equality regions.} For separation \(d\ge N^{-K_0}\), all selected losses are polynomial
in \(N\), and the unchanged pair factors retain quadratic suppression.
For \(d<N^{-K_0}\), use \(\theta_i=\mu+\rho\omega_i\), with
\(\sum\omega_i=0,\ \sum\omega_i^2=1\).
Permutation symmetry gives \(\partial_\rho G(\mu,0,\omega)=0\),
so \(-\partial_\rho G\ge c_R\rho/2\).
The \(Y\)-phase is exactly \(N\mu\); its real radial shifts do not
alter that phase. The coarea Jacobian in \(\mu,\rho\) is at least
\(c_RN\rho/2\). Convexity of the interval cube gives one interval
on each shape ray, and monotonicity gives multiplicity one there.

At a collision \(\theta_i=\theta_j\) the symmetric coupled shifts
are equal, so \(y_i-y_j\) vanishes to first order. The raw numerator
is bounded by \(C^{N^2}\rho^{N(N-1)}\), including its regular pair
denominators. After \(j\) Lie jets and one conservative flux/cutoff
loss, its remaining coarea power is
\(L=N^2-2j-4\ge4N_0+1>0\).
The exact inequalities \(d\le2\rho\) and \(\rho\le\sqrt N\,d\)
show that the near region
has \(\rho\le\sqrt N\,N^{-K_0}\), not literally \(\rho\le N^{-K_0}\).
Taking fixed \(K_0>1\) and increasing it if needed gives
\[
 C^{N^2}(\sqrt N\,N^{-K_0})^L
      e^{O_j(N\log(N+1))}(H+\log N)^2
 \le e^{C_jN\log(N+1)}(H+N)^{C_j}.
\]
The simple phase kernels have attenuation floors \(cN\eps\) in
this window. On each fixed shape ray, \(F=N\mu\) and bounded individual
phases put the unwrapped image inside a rectangle with side lengths
at most \(CN\). The multiplicity-one statement above, followed by
the interval estimate in Appendix~\ref{app:lie}, bounds its phase
integral by \(CN^2(H+\log N)^2\). This explicit polynomial factor
is included in \(e^{O_j(N\log(N+1))}\) in the displayed estimate.
The \(O(N^2)\) translated period boxes account for exactly this
factor. The unwrapped map still has multiplicity one.
On the separated region the coarea \(1/\rho\) costs at most a
polynomial.
Selected-diagonal fluxes vanish by the rational numerator, and
full-equality fluxes vanish by the positive collision exponent at
fixed \(\eps\). The actual Stokes multiplier is smooth with bounded
finite jets and its \(q\) remains on a plateau. This proves the
claim uniformly in \(\lambda\le\lambda_*\).
\end{proof}

\begin{lemma}[Protected large-current disks]\label{lem:protected}
On a named Stokes support, let \(\lambda_*\le\lambda\le1\).
For a sufficiently small constant \(c>0\), depending only on
the fixed point and previously chosen margins, the density is
holomorphic on
\[
 |t-t_\eps|\le\delta_t,\qquad \delta_t=c\lambda_*/N^2 .
\]
On this disk its product gaps are at least \(c'\lambda\),
its complete pairs have bound \(C^Ne^{-\kappa' N^2}\), and its
one-body factors have integrable majorants with total cost \(C^N\).
\end{lemma}
\begin{proof}
\noindent\textbf{Branch and compact-root control.} There is a named upper compact \(q\) with \(p_q=1\), so
\(P\ge1\). Fix a lower branch neighborhood inside \(m=1,p=0\)
and put \(a_B=\min|\sin\theta|>0\) there. With
\(\mu=\lambda P/N\), the exact imaginary-part formula gives
\[
 \Im W(t_\eps)\ge b_0\eps+\tfrac12\tau a_B\mu
 \ge b_\tau(\eps+\mu),\qquad
 b_\tau=\min(b_0,\tau a_B/2)>0.
\]
Indeed the increase from the outward shift is at least
\(\tau a_B\mu/2\), since \(\cosh v\ge1\) for real \(v\).
Write the disk constant as \(c_P\). After fixing \(\tau\), require
\[
 L_Sc_P\le b_\tau/4.
\]
Since \(\mu\ge\lambda_*/N\), throughout the disk
\[
 |\Delta W|\le L_Sc_P\lambda_*/N^2
 \le b_\tau \mu/(4N)\le b_\tau(\eps+\mu)/4.
\]
Thus these branch-near roots remain on the lower interior sheet.

Every non-plateau support angle has a fixed limiting distance
\(d_1>0\) from \(W=1\). Compute it on the preliminary endpoint
neighborhoods and the complement of the fixed lower plateau,
before shrinking \(\delta,\tau\); the upper branch is excluded.
All limiting
roots also have distance at least \(S_*\) from \(W=-1\).
For \(d=\min(d_1,S_*)>0\), the choices of
Lemma~\ref{lem:contraction}, followed by a reduction of \(c_P\),
ensure
\[
 \operatorname{dist}(W,\{1,-1\})
 \ge d-C(\eps+\tau+\delta)-L_S\delta_t\ge d/2.
\]
Here \(\delta\) is the fixed endpoint width.
Use the compact analytic continuations specified in that lemma.
Their motion is \(O(\delta_t)\); their modulus can exceed one.
The branch-coordinate motion is at most
\[
 C_\tau\frac{\delta_t}{\sqrt{\eps+\mu}}
 \le C_\tau c_P\frac{\sqrt{\lambda_*}}{N^{3/2}}.
\]
Reduce \(c_P\) further so these motions consume less than
the reserved \(q'-q_1\) and \(\eta/2\) pair margins and retain
the compact denominator margins. The regular double-branch
expression controls the remaining neighborhood. This proves
same-group bound \(q'\), cross bound \(1+\eta\), and the named
bound on the whole disk without a unit-disk Schur argument there.
The direct count in Lemma~\ref{lem:contraction} gives the same
quadratic exponent, including deletion of at most \(j\) pairs.

\par\smallskip\noindent\textbf{Product gaps and one-body integration.} At the radial center the named upper root has
\(|z_q|\le e^{-c_q\lambda}\), \(c_q>0\), because its inward
shift is \(2\tau\lambda\) and its arc has positive sine.
Branch roots remain interior on the disk and compact roots obey
\(|z_i(t)|\le |z_i(t_\eps)|e^{C\delta_t}\).
Adding \(Cc_P\le c_q/2\) to the disk restrictions gives
\[
 |Z|\le e^{-c_q\lambda+CN\delta_t}\le e^{-c_q\lambda/2}.
\]
Also \(|Y|\le e^{-3\tau\lambda P/2}\), independently of \(t\).
Both product gaps are therefore at least \(c'\lambda\).
Choose \(c_P\) after \(\tau\) and before the final \(\eps_0\).
This order keeps the disk within the prescribed branch sheet and
establishes the geometric bounds before estimating the current integral.

Finally put \(\mu=\lambda P/N\). The real displacement of the
branch center is \(O(\eps^2+\mu^2+\delta_t)\), and the imaginary
margin is at least \(c_2(\eps+\mu)\).
Split into \(|u|\le C(\eps+\mu)\) and its complement.
The imaginary margin bounds the square root in the first
case, and the linear real term bounds it in the second, with
the earlier choices of \(\tau\) and then \(c_P\) sufficiently small. Hence
\[
 |R(y;t)|\le C(|u|+\eps)^{-1/2}.
\]
This is a majorant independent of the coupled scalar \(P\);
no integration treats \(P\) as a free constant. Compact
one-body factors are bounded. The inverse-product numerator,
the exact Jacobian, and its Stokes replacement minor cost
\(C^N\). These facts prove all the stated assertions.
\end{proof}

\begin{proposition}[Large-\(\lambda\) current]\label{prop:largeS}
For all odd $N>N_0$, the differentiated current integrated over
\(\lambda_*\le\lambda\le1\) satisfies \eqref{eq:reqLS}.
\end{proposition}
\begin{proof}
Lemma~\ref{lem:protected} bounds the absolute integral of
the undifferentiated density on its common disk by
\(C^Ne^{-\kappa N^2}\lambda^{-2}\), including the named-index
count in \(C^N\). The majorant is integrable and uniform on
that disk, so the integrated density is holomorphic there.
Cauchy's inequality of order \(j\) adds
\(C_jN^{2j}\lambda_*^{-j}\). Integration in \(\lambda\) costs
at most \(\lambda_*^{-1}\). Summing the convergent series
\(\sum_N C^NN^{2j}e^{-\kappa N^2}\) gives even the stronger
bound \(C_j\lambda_*^{-j-1}\). Its weaker version
\(C_j\lambda_*^{-j-2}\) is \eqref{eq:reqLS}.
By \eqref{eq:alphachoice}, this is \(o(\eps^{-1/2})\).
\end{proof}

\begin{proposition}[The remaining high orders]\label{prop:highKS}
The bound \eqref{eq:reqhigh} holds on
\(N\ge D\sqrt H\).
\end{proposition}
\begin{proof}
Use Lemma~\ref{lem:originaldisk}, which establishes the original
\(c\eps\) disk, the \(c\eps\) product gaps, integrable one-body
majorants, and complete pair bound \(C^Ne^{-\kappa N^2}\).
This invocation is independent of the large-current disk.
Cauchy of order \(j\) therefore gives
\(C_j^N\eps^{-j-2}e^{-\kappa N^2}\).
Integrating the small-\(\lambda\) current only reduces this bound.
The argument needs no enlarged Cauchy disk and makes no
unbounded-\(N\) assertion that \(1-r^N\ge cN\eps\).
\end{proof}




\section{The onsite series and the infinite odd-particle tail}\label{sec:conclusion}

The onsite term is part of each full coefficient and is estimated
separately before the tail is assembled. Its contour has no global
product poles, so a one-phase regularization suffices. The final theorem
then combines the onsite bound with the four offsite estimates from
Section~\ref{sec:tailreduction}.

\begin{proposition}[The onsite differentiated series]\label{prop:onsitesum}
At either selected boundary half, for every fixed nonnegative integer
$j$, there are constants such that, with $H=\log(1/\eps)$,
\[
 \sum_{\substack{N\ge1\\N\ {\rm odd}}}|f_N^{(j)}(t_\eps)|
 \le C_j\exp\bigl(B_j\log^2(H+1)\bigr).
\]
In particular this sum is $o(\eps^{-\delta})$ for every $\delta>0$.
\end{proposition}
\begin{proof}
First work in the positive-$S_*$ geometry. The reduced onsite form
has global factor $(YZ)^{-1}$ and the same complete pairs and
one-body measures. The same homotopy and weighted Cartan identity
give $f_N=F_{00,N}+K_{00,N}+S_{00,N}$, with precisely the same
current multiplier $-2\ii\tau\,\partial_q\psi$. There are no
global product poles. The proof of Proposition~\ref{prop:F} applies
with these denominators deleted, and its all-$N$ summable majorant
therefore bounds $\sum|F_{00,N}^{(j)}|$ by the displayed subpower
function of $\eps^{-1}$.

For $K_{00,N}$ and each named current use a fixed product partition
into a true lower-branch plateau and its regular compact complement.
All branch cutoff supports and their jets are inside $m=1,p=0$;
the named current index is compact. Replace just the branch angles
by their regular coordinates $\phi_i=\arccos W_i$. The branch
one-body form and complete pairs are analytic in these coordinates,
including at collisions, as in \eqref{eq:branchregular}. Set
\[
 V_i=\frac{\partial_t\phi_i}{\partial_{\theta_i}\phi_i}
 \quad\text{on branch indices},\qquad V_i=0
 \quad\text{on compact indices}.
\]
The quotient is the analytic bounded function $S'/W_{\theta_i}$.
This field freezes each branch phase. On the actual branch plateaus
$p'_i=m'_i=0$, so $V\cdot\nabla P=0$ and
$V\cdot\nabla\sum m_i=0$, including all parameter derivatives.
Thus the residual field on all hybrid coordinates is zero. The
pullback recurrence \eqref{eq:pullbackLie} has only regular
coefficient differentiation and derivatives of the fixed real
cutoffs; it has no selected-difference pole.

At each fixed derivative order a product-rule descendant changes
at most $j$ complete pairs. The unchanged factors retain
$C_j^N e^{-\kappa N^2}$ by Lemma~\ref{lem:contraction}. The
analytic one-body and pair jets, inverse products, exact coupled
Jacobian, and cutoff derivatives cost $C_j^N N^{C_j}$. This holds
uniformly in all $N$ and $0\le\lambda\le1$, after the fixed
branch margins and deformation strength are selected. When pulled
back to real angles, the regular branch measures have uniformly
bounded integrals: the estimate
$|\phi_{\theta_i}|\le C(|u_i|+\eps)^{-1/2}$ from the original
disk proof is independent of the coupled occupancy. It remains
valid for the current on the full homotopy, with the fixed small
$\tau$ and the imaginary margin $\eps+\lambda P/N$.

There are no puncture boundaries or collision poles in this
one-phase regularization. Integration by parts on the real torus
is justified by smooth cutoffs at fixed $\eps,N,\lambda$; the
regular branch coordinate cancels its one-body singularity.
Consequently
\[
 |K_{00,N}^{(j)}(t_\eps)|+
 |\mathcal J_{00,N,j}(\lambda)|
 \le C_j^N N^{C_j}e^{-\kappa N^2}.
\]
Integrate over $[0,1]$ and sum over all odd $N$. The series is
convergent, which proves the assertion together with the $F$ bound.
For the negative family, \eqref{eq:signidentity} and conjugation
give the same absolute estimate. The estimate follows directly from the individual onsite integrals,
so it also controls the sum of their absolute derivatives.
\end{proof}

\begin{remark}[Independent physical check of the onsite sum]
The classical per-site identity in \cite[eqs.~(4.1)--(4.2)]{Orrick2001}
also gives an independent check of the full-site normalization.
At the origin its angular integrand has no displacement multiplier:
by Lemma~\ref{lem:residue} it is precisely $f_N$.
Since $\langle\sigma_{00}^2\rangle=1$, for real $t>1$ it follows that
\[
 t(1-t^{-4})^{1/4}\sum_{N\ {\rm odd}}f_N(t)=1.
\]
Local normal convergence and the identity theorem extend this equality
throughout $|t|>1$, with the fourth root inherited from infinity.
Consequently
\[
 \mathcal X_H(1/t)=1+2t(1-t^{-4})^{1/4}
                    \sum_{N\ {\rm odd}}T_N(t).
\]
This supplies a second final bridge using only the offsite tail.
The preceding proposition proves the stronger absolute onsite
derivative bound directly and also permits the full-coefficient
tail statement below. The check uses the physical per-site identity and the identity theorem
on the exterior domain.
\end{remark}

\begin{theorem}[Uniform differentiated full bulk tail]\label{thm:tail}
At every selected $t_*$ and for every $0\le j\le k$,
\[
 \sum_{\substack{N>p\\N\ {\rm odd}}}|T_N^{(j)}(t_\eps)|
 =o(\eps^{-1/2}),\qquad
 \sum_{\substack{N>p\\N\ {\rm odd}}}|C_N^{(j)}(t_\eps)|
 =o(\eps^{-1/2}).
\]
Here the first sum denotes the ordinary unsigned offsite coefficients;
the signed estimates are used to cover the negative family.
\end{theorem}
\begin{proof}
\noindent\textbf{Differentiate before splitting.} Differentiate the locally normally convergent original series first,
as justified in Lemma~\ref{lem:schur}. At each $N$, use the exact
$F+K+S$ identity. All moving particle-number and current-window
splits below are introduced only after differentiation.

\par\smallskip\noindent\textbf{The selected contour, large current, and high orders.} The $F$ sum over all $N>p$ is at most
$\exp(O_j(\log^2(H+1)))$. The large-$\lambda$ current over all
$N>p$ is at most $C_j\eps^{-\alpha_0(j+2)}=o(e^{H/2})$,
because $\alpha_0<1/[2(k+2)]$.
Choose a common positive $\kappa$ smaller than the fixed pair
exponents, and then choose $D$ with $\kappa D^2>k+2$.
The sum of the original and small-current contributions for
$N\ge D\sqrt H$ is bounded by
\[
 \sum_{N\ge D\sqrt H}C_j^N\eps^{-j-2}e^{-\kappa N^2}
 \le \exp\bigl((j+2-\kappa D^2)H+O_j(\sqrt H)\bigr)
 =o(e^{H/2}).
\]
This estimate uses the whole remaining infinite tail, with no
upper particle-number cutoff.

\par\smallskip\noindent\textbf{The intermediate window and full-site conclusion.} Every intermediate order $p+2\le N<D\sqrt H$ is covered by the
fixed partition and Propositions~\ref{prop:micro}, \ref{prop:allB},
\ref{prop:mixed}, and \ref{prop:compactR}. The positive collision
power is
\[
 N^2-2j-4\ge(p+2)^2-(p^2-1)-4=4p+1>0.
\]
Their summed bound has logarithm
$O_j(\sqrt H\log(H+1))=o(H)$. The small-current interval has
length at most one. Combining the pieces proves the first estimate
for either signed kernel at positive $S_*$. Negation and
conjugation give the original coefficient on the negative family.
Finally \eqref{eq:onsiteidentity} and
Proposition~\ref{prop:onsitesum} prove the full-site tail estimate.
\end{proof}

\section{Noncancellation and the natural boundary}

We now compare the first singular coefficient with the complete tail
at a fixed selected point. The nonzero leading derivative survives both
physical normalization factors. Density then converts these local
obstructions into nonextension at every point of the unit circle.
\begin{proof}[Proof of Theorem~\ref{thm:main}]
Put $U(t)=\sum_{N\ge1,\,N\ {\rm odd}}C_N(t)$. For each selected
point, the finitely many lower terms have bounded derivatives.
Theorems~\ref{thm:first} and \ref{thm:tail} give
\[
 U^{(k)}(t_\eps)=4L_*\eps^{-1}+o(\eps^{-1}),
 \qquad U^{(j)}(t_\eps)=o(\eps^{-1})\quad(0\le j<k).
\]
Define $\widehat{\mathcal X}(t)=\mathcal X_H(1/t)$.
By \eqref{eq:bulk},
\[
 \widehat{\mathcal X}(t)=Q_H(t)U(t),\qquad
 Q_H(t)=t(1-t^{-4})^{1/4}.
\]
At our selected nonaxis points $Q_H$ is locally analytic and
nonzero on its chosen continuation. Leibniz's rule therefore gives
\[
 \widehat{\mathcal X}^{(k)}(t_\eps)
 =4Q_H(t_*)L_*\eps^{-1}+o(\eps^{-1}),
 \qquad Q_H(t_*)L_*\ne0.
\]
A holomorphic continuation through $t_*$ would bound this finite
derivative on a short exterior radial segment, a contradiction.
The selected points are dense in both upper quadrants, and
conjugation supplies both lower quadrants. Every neighborhood of
every point of $|t|=1$, including an axis point, contains one of
these obstructing points. Thus no point of the circle permits
local holomorphic continuation of the exterior germ.

The inversion $s=1/t$ is biholomorphic in a neighborhood of each
unit-circle point and transfers the conclusion to the prescribed
high-temperature germ. Finally
$\beta_{\rm phys}(s)=\operatorname{arsinh}(s)/(2J)$ is locally
analytic and nonzero at each selected nonaxis point on the
physical continuation. Multiplying by it preserves the nonzero
leading derivative and the dense obstruction. This proves the
statement for the physical susceptibility $\chi_H$.
\end{proof}

\paragraph{Quantifiers and scope.}
For each selected point the prime $p$, derivative order $k$, all
cutoffs, and all constants are fixed before $\eps\to0$. Uniformity
in $p$ is neither used nor asserted. Uniformity in particle number
is exactly the one established in the signed tail and onsite
estimates. The result concerns the complete bulk sum, with no
particle truncation and no restriction to a lattice diagonal.

\appendix
\section{Fixed-cutoff differentiation, pole counting, and coarea}
\label{app:lie}

This appendix provides the calculus used in the intermediate-order
sectors. The order of operations is fixed: differentiate on a smooth,
fixed real domain; verify that the puncture fluxes vanish; then estimate
the resulting absolute integral by coarea. Freezing the two product
phases keeps their kernels simple during differentiation.

\subsection{The identity on a fixed real integration domain}
\begin{lemma}[Fixed-domain differentiation]\label{lem:lie}
Let \(w\) be a smooth, \(t\)-independent weight on a real angular torus.
For a \(t\)-dependent complex vector field \(V\) and a top form \(A_t\),
set \(\mathcal D=\partial_t-\Li_V\). If all puncture fluxes vanish
at every intermediate stage through order \(j\), then
\[
 \partial_t^j\int wA_t=\int\mathcal D^j(wA_t).
\]
If \((\partial_t-V\cdot\nabla)Y
=(\partial_t-V\cdot\nabla)Z=0\), the factors
\((1-\sigma Y)^{-1}\) and \((1-\sigma Z)^{-1}\) remain simple throughout this identity.
\end{lemma}
\begin{proof}
On a top form, \(\Li_V=\dd\iota_V\) is a divergence. Its integral is the
boundary flux, including every puncture boundary. Its vanishing proves
the first derivative identity. Apply that identity to the entire
\(t\)-dependent top form produced at the preceding step. In particular,
the iteration differentiates \(V\), its divergence, and all coefficients
already present. The kernel assertion follows from the derivation rule.
Complex coefficients in \(V\) are allowed by linearity of real
integration by parts. No real flow or holomorphic extension of \(w\)
is required.
\end{proof}

If \(\sum_h w_h=1\) is a fixed smooth partition, use a different field
\(V_h\) on each weighted term:
\begin{equation}
 \partial_t^j\int A_t
   =\sum_h\int(\partial_t-\Li_{V_h})^j(w_hA_t).
 \label{eq:chartwiseLie}
\end{equation}
Apply the preceding lemma to each weighted integral and add the
results. Each field is used only for its own integral; its two
frozen-product identities are preserved independently of the fields
on overlapping supports.

\subsection{Regular branch coordinates and the selected-pair field}
On the original all-branch chart let \(\chi_t:u\mapsto\phi\), where
\(\phi=\arccos W\). The inverse \(y_t(\phi)\) is even and analytic
on a fixed polydisk about \((t_*,0)\), because \(W_\theta\ne0\) there.
The one-body form and the complete pair are
\begin{equation}
 \frac{\dd y}{2\pi\ii}\frac{z}{\sinh\gamma}
   =-\frac{z}{\pi(1-y^{-2})}\dd\phi,\qquad
 P_{il}=(\phi_i-\phi_l)^2 A_t(\phi_i,\phi_l).
 \label{eq:branchregular}
\end{equation}
Here \(A_t\) is analytic with bounded fixed jets. Indeed,
\(y_t(\phi_i)-y_t(\phi_l)\) has the two factors
\(\phi_i-\phi_l\) and \(\phi_i+\phi_l\); the latter cancels the
corresponding factor in \(1-e^{-\ii(\phi_i+\phi_l)}\).
No division by a remaining collision zero is used.

Write the regular top form as
\[
 \Omega_t=
  \frac{\mathscr N_t(\phi)}{\mathcal K_t(\phi)}
       \dd\phi_1\cdots\dd\phi_N,\qquad
 \mathcal K_t=(1-\sigma Y_t)(1-\sigma Z),\qquad
 \mathscr N_t=\Delta(\phi)^2\mathscr A_t(\phi).
\]
The analytic factor \(\mathscr A_t\) includes all normalized measures,
the inverse-product numerator, and the regular pair coefficients.

Put \(a_i=S'/W_{u_i}\), \(b_i=\partial_{u_i}\phi_i\), and
\(A=\sum_i a_i\). For a selected pair \(p,q\), set
\begin{equation}
 \begin{split}
 V_i&=a_i\quad(i\ne p,q),\\
 V_p&=a_p+\frac{Ab_q}{b_p-b_q},\qquad
 V_q=a_q-\frac{Ab_p}{b_p-b_q}.
 \end{split}\label{eq:twofield}
\end{equation}
Then \(\sum_iV_i=0\) and
\(\sum_i b_iV_i=\partial_t\sum_i\phi_i\), so \(Y,Z\) are frozen.
In regular coordinates the residual field
\(B=\partial_t\chi_t-(\chi_t)_*V\) has only two nonzero entries:
\[
 B_p=-\frac{Ab_pb_q}{b_p-b_q},\qquad B_q=-B_p.
\]
Since \(b(\phi)=B_t(\phi)/\phi\) with \(B_t\) even, analytic, and
nonzero near zero,
\[
 \frac{b_pb_q}{b_p-b_q}
     =\frac{H_t(\phi_p,\phi_q)}{\phi_q-\phi_p}.
\]
The analytic coefficients and any fixed number of their derivatives
have polynomial \(N\)-bounds. Both \(V\), expressed in regular
coordinates, and \(B\) have at most one selected-difference pole;
\(\operatorname{div}_\phi B\) has at most two.

The pullback identity needed for smooth real cutoffs is
\begin{equation}
 (\partial_t-\Li_V)(w\chi_t^*\Omega_t)
  =w\chi_t^*(\partial_t+\Li_B)\Omega_t
       -(V\cdot\nabla_u w)\chi_t^*\Omega_t.
 \label{eq:pullbackLie}
\end{equation}
For clarity, if \(J_t=\det D_u\chi_t\), then
\(\partial_tJ_t=J_t(\operatorname{div}_\phi
(\partial_t\chi_t\circ\chi_t^{-1}))\circ\chi_t\).
The corresponding density formula for \(\mathcal L_V\) is
\(V\cdot\nabla_u+\operatorname{div}_uV\).
Subtracting these two product rules yields \eqref{eq:pullbackLie},
including the displayed derivative of the real weight. In particular,
\((\partial_t+B\cdot\nabla_\phi)\mathcal K_t=0\).
Formula \eqref{eq:pullbackLie} is used on the real \(u\)-domain;
the smooth function \(w\) is never extended holomorphically in \(\phi\).

\subsection{The complete high-order recurrence}
\begin{lemma}[Unfactored jet filtration]\label{lem:jets}
For each fixed \(j\), the \(j\)-fold operation in
\eqref{eq:pullbackLie} is a finite sum of terms
\begin{equation}
 (\partial_u^\mu w)(u)\,
 \chi_t^*\left[
  \frac{c_{j,\mu,\ell,\nu}(t,\phi)\,
        \partial_t^\ell\partial_\phi^\nu\mathscr N_t(\phi)}
       {\mathcal K_t(\phi)}
       \dd\phi_1\cdots\dd\phi_N\right].
 \label{eq:jetexpansion}
\end{equation}
If \(m\) is the selected-difference pole order of a displayed
coefficient, then that term can be chosen to satisfy
\begin{equation}
 m+|\mu|+|\nu|\le2j,\qquad
 \ell+|\mu|+|\nu|\le j .
 \label{eq:jetdegree}
\end{equation}
Apart from these explicitly counted poles, fixed coefficient jets
are bounded by \(C_jN^{C_j}\) on the fixed regular-coordinate
neighborhood. There are at most \(C_jN^{C_j}\) terms.
\end{lemma}
\begin{proof}
For \(j=0\) take \(c=1\). Suppose \eqref{eq:jetexpansion} holds.
Apply \eqref{eq:pullbackLie} to each term, and use the fact that
the kernel is frozen. The possibilities are as follows:
\begin{enumerate}
\item A fixed-\(\phi\) \(t\)-derivative on \(c\) does not raise its
difference-pole order: \(\delta=\phi_p-\phi_q\) is independent of
\(t\) in these coordinates.
\item \(B\cdot\nabla_\phi c\), or multiplication by
\(\operatorname{div}_\phi B\), raises \(m\) by at most two.
\item \(B\cdot\nabla_\phi\) on the numerator derivative raises
\(|\nu|\) by one and \(m\) by at most one.
\item \(\partial_t\) on the numerator derivative raises \(\ell\)
by one without adding a difference pole.
\item The final term of \eqref{eq:pullbackLie} raises \(|\mu|\)
by one and contributes one component of \(V\), whose pole
order is at most one.
\end{enumerate}
Every operation therefore raises the first count by at most two
and the second count by at most one. These are all product-rule
terms, including derivatives of the previous field coefficients.
Bounded analytic coefficients, at most \(N\) summands in \(A\),
and a fixed number \(j\) of differentiations give the claimed
polynomial term and coefficient bounds.
\end{proof}

The object differentiated in this recurrence is the
\emph{unfactored numerator} \(\mathscr N_t\).
A numerator derivative of order \(\ell+|\nu|\le j\) changes at
most \(j\) pair factors. The remaining pair contractions survive.
Fixed-\(\phi\) \(t\)-derivatives preserve the Vandermonde polynomial;
\(|\nu|\) coordinate derivatives reduce its full-collision degree
by at most \(|\nu|\).
Keeping the Vandermonde inside the differentiated numerator preserves
the stated selected-pole count. Dividing by it would introduce additional
poles at unselected collisions, so Lemma~\ref{lem:jets} is applied
to the unfactored numerator throughout.

\subsection{Rational weights and all boundary fluxes}
Fix an integer \(M>2k+2\), independent of \(N\), and put
\[
 d=\max_{i<l}|u_i-u_l|,\qquad
 w_{pq}=\frac{(u_p-u_q)^{2M}}
                  {\sum_{i<l}(u_i-u_l)^{2M}}\quad(d>0).
\]
They sum to one off full equality. For \(|\mu|\le k\),
\begin{equation}
 |\partial_u^\mu w_{pq}|
 \le C_kN^{C_k}
       \frac{|u_p-u_q|^{2M-|\mu|}}{d^{2M}}.
 \label{eq:weightjets}
\end{equation}
To see this, differentiate the quotient. A derivative of the
denominator contributes a power at most \(d^{-1}\); replace it,
when needed, by \(|u_p-u_q|^{-1}\). The denominator itself is
at least \(d^{2M}\), and its derivative sums have polynomial
\(N\)-cost. This proves \eqref{eq:weightjets} term by term.

On the original branch cone,
\(|\phi_p-\phi_q|\ge c|u_p-u_q|\), by
Lemma~\ref{lem:branch}. Combining \eqref{eq:weightjets} with
\eqref{eq:jetdegree} absorbs every selected pole at a cost
\(C_jN^{C_j}d^{-2j}\) on a separated support, while the
undifferentiated pair factors retain their quadratic contraction.
Other fixed smooth cutoffs are included by the ordinary product
rule. A cutoff at scale \(b_N\) or \(\rho_N\) has its corresponding
inverse-scale cost; these costs are retained in
Propositions~\ref{prop:micro} and \ref{prop:allB}.

Near full equality, the numerator has degree \(N(N-1)\).
The bound \(m+|\mu|+|\nu|\le2j\) shows that a full \(j\)-th
density loses at most \(2j\) powers of the shape radius.
A further inverse radius is reserved when a boundary flux or
a transition-cutoff estimate requires it. This is the origin
of the conservative loss \(2j+1\) used in the sector proofs.

The same power counts verify the flux hypothesis of
Lemma~\ref{lem:lie} at every preceding differentiation stage.
First take a selected-diagonal tube of radius \(h\) on a region
where \(d\) is bounded below. The pair double zero, the
\(2M\)-power weight, the preceding \(j\) operations, and one
extra field factor give a flux at most
\[
 C_{\eps,N,j}\,h^{2M+1-2j}.
\]
It vanishes. Next take a full-equality sphere of radius \(\rho\).
At fixed \(\eps>0\) all global denominators and regular coordinate
changes are bounded; the shape surface measure contributes
\(\rho^{N-2}\). A sufficient flux exponent is
\[
 N(N-1)-2j-1+N-2=N^2-2j-3>0
 \quad(N\ge N_0+2,\ j\le k).
\]
The limits are taken in this order at fixed \(\eps,N\).
There are no boundaries at a cutoff support endpoint because
the cutoffs are smooth on the closed torus. These flux checks
apply at every earlier stage \(0\le l\le j\), and establish the
iterated identity before any boundary-uniform estimate is taken.

\subsection{Mixed and compact charts}
In a mixed chart only the true branch coordinates are replaced
by \(\phi_i\). The selected denominator is \(1-b_q/b_j\), bounded
away from zero on the nested support, and \(1/b_j\) is analytic
at the branch. There is consequently no selected-collision pole
in this calculation. The same pullback recurrence includes the
regular branch one-body form, the compact-angle measures, and
the full contour Jacobian. All fixed jets have polynomial
\(N\)-cost, apart from \(C^N\) one-body products.

For the compact right chart the coefficients are smooth in real
angles and analytic in \(t\). The factorization
\(D_{pq}=(\theta_p-\theta_q)C_{pq}\), with \(|C_{pq}|\ge c>0\),
uses smooth division, as proved in Proposition~\ref{prop:compactR}.
The same five product-rule cases prove a two-power loss per
operation in the real angular variables. There are no singular
branch-coordinate changes in this compact case. It does not
require analytic continuation of a deformation bump.

\paragraph{Coupled occupancy and changed pairs.}
The compact-right current field may move a coordinate where
\(p'\ne0\). The exact contour Jacobian gives
\[
 (V\cdot\nabla\log y_i)_{\rm occ}
   =\frac{\lambda\tau}{2N}m_i\sum_l V_l p'_l.
\]
An individual off-diagonal entry is \(O(\lambda\tau/N)\).
The compact-right field has two nonzero components, each possibly
\(O(N/|\theta_p-\theta_q|)\), so the contraction can be
\(O(\lambda\tau/|\theta_p-\theta_q|)\). It moves every
\(m\)-supported root and must be retained. All these roots are
compact in this chart; no singular raw branch-angle jet is used.

To count descendants, keep each complete pair \(F_e=P_{il}\)
as a scalar function of all angles and \(t\), including its
occupancy dependence. Write the density as \(G\prod_e F_e\).
The scalar operator \(\partial_t-V\cdot\nabla\) is a derivation;
the top-form correction is multiplication by
\(-\operatorname{div}V\). In one product-rule summand a step
hits a coefficient or \(G\), an already changed pair, or one
previously unchanged pair. Divergence changes none. Induction
therefore leaves at most \(j\) changed pairs per descendant at
order \(j\), even though the sum of descendants involves all pairs.
In regular coordinates, hitting either the double-zero polynomial
or its analytic coefficient marks that same pair. No division
by an original pair is involved.

There are \(O(N^2)\) possible pair choices at a step. Expanding a
coordinate contraction or an occupancy sum introduces at most
\(O(N)\) further choices. A fixed number of chain rules therefore
has \(N^{O(j)}\) descendants. Fixed jets of \(P/N\) are bounded
by the fixed bump jets, and fixed compact root and pair jets are
bounded on the continuation neighborhoods of
Lemma~\ref{lem:contraction}. Sums of phase jets, and inverses of
the separated smooth divided determinant \(C_{pq}\), have only
polynomial fixed-jet costs. The explicit selected difference pole
still adds at most two powers per Lie step. One-body and Jacobian
products cost \(C_j^N N^{O(j)}\). Consequently the coupled terms
fit the existing polynomial coefficient and term budgets and
the unchanged pairs retain \(C_j^N e^{-\kappa N^2}\).
In mixed Stokes charts the plateau identities
\(V\cdot\nabla P=V\cdot\nabla\sum m_i=0\) remove this coupling;
original \(K\) has \(\lambda=0\).

\subsection{Simple kernels and the order of coarea}
For either signed kernel, translate its phase by zero or $\pi$. For \(a>0\),
\[
 \frac1{|1-e^{-a+\ii \mu}|}
 \le\frac C{\min(1,a)+\operatorname{dist}(\mu,2\pi\Z)}
 \quad\bigl(0\le\operatorname{dist}(\mu,2\pi\Z)\le\pi\bigr).
\]
On an interval of length \(L\), the integral is at most
\(C(1+L)(1+|\log a|)\) for \(0<a<1\). A variable damping can be
replaced by its positive lower bound in this estimate.

For two unwrapped phases \(F,G\), suppose their image is contained
in \(I_F\times I_G\), the map has multiplicity at most \(m_0\),
and its real Jacobian has already been divided out in the density
bound. If the attenuation floors are \(a_Y,a_Z\in(0,1)\), the
remaining two-kernel integral is bounded by
\[
 C m_0(1+|I_F|)(1+|I_G|)
       (1+|\log a_Y|)(1+|\log a_Z|).
\]
Indeed extend the nonnegative integrand to the containing rectangle
and apply the one-dimensional bound twice. Equivalently, subdivide
each interval into translates of one period. This accounts for
periodic wrapping separately from the multiplicity of the unwrapped
map. In the near-equality sectors \(m_0=1\) on each fixed shape
ray and both lengths are at most \(CN\), yielding an explicit
\(CN^2\) factor. Any already counted index cells or fixed-jet
terms retain their stated polynomial costs. Such factors are
absorbed by the existing \(e^{C_jN\log(N+1)}\) budgets.

Coarea is applied only after all differentiation, to maps whose
real Jacobian and multiplicity are specified in the sector proof.
Sign and extreme-index cells are introduced only in that final
absolute integral, after the smooth fixed-domain identity
\eqref{eq:chartwiseLie} has been established. When a shape radius
is exchanged for a phase variable, first bound its remaining positive
power on the support and then integrate the phase. The exchanged
radius is already included in this change of variables.




\section{Constant and quantifier dependencies}\label{app:constants}
The following order makes every uniformity statement in the tail
proof explicit. Each choice uses only earlier quantities; all choices
precede the final restriction on \(\eps\) and are independent
of the particle number.
\begin{enumerate}
\item Fix \((p,a,b)\), hence \(N_0,k,S_*,\theta_B\), and the
original radial constant and a preliminary bound on \(|S'|\).
The final disk constants are chosen in step 4. There are only finitely many
derivative orders \(0\le j\le k\).
\item Fix a sufficiently small branch neighborhood and compact
margins. Choose the nested ratio \(L\) large enough for the
outer-anchor/true-branch slope separation, also relative to the
bounded upper-current slopes on a preliminary endpoint neighborhood
separated from the upper branch. Fix the smooth nested
partition and all of its finite derivative bounds.
\item Obtain the same-group pair constant \(q<1\) at the limiting
real support, with the groups and continuation neighborhoods of
Lemma~\ref{lem:contraction}. Choose \(q<q_1<q'<1\), then the cross-group slack
\(\eta\), and then the upper endpoint slivers and deformation
strength \(\tau\), retaining the reserved same/cross disk slack.
The upper branch must lie inside \(a=p=1\). Fix the bump jets
and the positive constants \(b_\tau,c_q\) at this step.
\item Choose the selected \(c/N\), original \(c\eps\), and
protected-current \(c\lambda_*/N^2\) disk constants small enough
for the already fixed margins. In particular impose
\(L_Sc_O\le b_0/4\), \(L_Sc_P\le b_\tau/4\), the
compact-distance and reserved pair-slack restrictions, and
\(Cc_P\le c_q/2\) from Lemma~\ref{lem:protected}.
No reduction of \(\tau\) follows this choice. Take a common smaller positive
\(\kappa\) for the finitely many pair estimates.
\item Choose \(0<\alpha_0<1/[2(k+2)]\), and then \(D\) with
\(\kappa D^2>k+2\). Fix the rational-weight integer \(M>2k+2\)
and a sufficiently large compact-equality exponent \(K_0>1\).
\item Choose the resultant exponent \(A\), followed by the small
microcore prefactor \(c_*\) in \(b_N=c_*e^{-(A+4)N}/N\).
Determine the finite-jet constants on the anchor neighborhood
\(\rho\le b_N/8\), \(|u_a|\ge b_N/2\). Its lower bound on
\(|\mu|\) determines the cubic coefficient \(C\) independently of
the still unselected exponent \(B\).
\item With the finite \(\Gamma_*\) defined in
Proposition~\ref{prop:allB}, choose
\(B,C_1>A+4+\max\{1,\Gamma_*\}\). This ensures the geometric
inclusions \(2\rho_N\le b_N/8\) and \(\delta_N\le b_N/16\)
for every \(N\ge N_0+2\), including the finitely many low orders.
Increase \(B\) further so that
\(\exp((A+4)N^3/2+\tfrac12N^2\log N+C_jN^2)
 (2e^{-BN})^{N^2-2j-4}\) is controlled for all large
\(N\) and every \(j\le k\). Only the remaining numerical bounds
at finitely many orders are absorbed into a point-dependent
constant; their support geometry already satisfies the inclusions.
\item Finally decrease a single \(\eps_0>0\) for all the conditions.
In the intermediate window, for every fixed \(C,m,r\),
\[
 \eps^r e^{CN}N^m\longrightarrow0\quad(r>0),\qquad
 \eps/b_N^m\longrightarrow0,
 \quad N\le D\sqrt{\log(1/\eps)}.
\]
These limits also give \(N^{C_j}(\eps+\lambda_*)\to0\).
They are used only in that window.
In particular impose \(N\eps e^{AN}\le c\),
\(\eps/b_N\le c\), and the finitely many compact-jet
perturbation inequalities for \(j\le k\), as well as
\(\eps+\lambda_*\) sufficiently small relative to the fixed
branch and compact margins. For the first two the logarithms
at \(N\le D\sqrt H\) are bounded by
\(-H+O(\sqrt H+\log H)\); for the compact perturbations the bound
is \(-\alpha_0H+O(\log H)\). Thus one \(\eps_0\), independent
of \(N\), enforces every required condition. Only finitely many
fixed exponents from the preceding estimates are needed.
\end{enumerate}

The distinction between costs is essential: an ordinary one-body
product costs \(C_j^N\); a full collision coefficient may cost
\(e^{C_jN^2}\); conversion from the exponentially small branch
anchor to shape coordinates may cost \(e^{CN^3}\). The last cost
is paid only where the positive collision power at \(\rho_N\)
absorbs it. It is not replaced by a polynomial in \(N\).

\begin{longtable}{>{\raggedright\arraybackslash}p{.22\textwidth}>{\raggedright\arraybackslash}p{.30\textwidth}>{\raggedright\arraybackslash}p{.38\textwidth}}
\toprule
Quantity & May depend on & Uniformity actually used\\
\midrule
\endhead
\(c_0,c_1\), original radius/disk &
Fixed selected point & All original angles and all odd \(N\);
radius held fixed under each \(t\)-derivative.\\
Branch and compact margins &
\(S_*,\theta_B,k\) & Independent of \(N,\eps\); compact groups
exclude the exact branch before strict constants are chosen.\\
\(q,\eta,\kappa,\tau\) &
The preceding margins and point & Entire stated support, all
occupancies \(P/N\in[0,1]\); radial centers and both stated disks.\\
Selected-contour disk \(c/N\) &
Point, margins, \(\tau\) & All \(\eps\) sufficiently small;
not an assertion about the original undeformed contour on that disk.\\
Current disk \(c\lambda_*/N^2\) &
Point, margins, \(\tau\) & \(\lambda_*\le\lambda\le1\);
compact roots may leave the unit disk.\\
\(C_j\) in jets &
Point, fixed bumps, \(j\le k\) & Independent of \(N,\eps\) within
the stated window; powers \(N^{C_j}\) are explicit.\\
Final \(C_j\) in \eqref{eq:reqinter} &
Point, \(j\), all preceding margins, \(\tau,\kappa,\alpha_0,D\),
\(M,K_0,A,c_*,B,C_1\) & Independent of \(N,\eps,\lambda\).
Finite small-order constants may be large but are
\(\eps\)-independent.\\
\(A,b_N,\rho_N\) &
Algebraic point; \(B\) also on \(k\) & Exponential \(N\)-dependence
is retained. \(\eps e^{cN}\to0\) only for \(N\le D\sqrt H\).\\
\(\alpha_0,D\) &
Point and \(k\) & Fixed before \(\eps\to0\); all lower derivatives
use the same finite choices.\\
First-chart constants &
The fixed \(N_0,p,a,b\) & No uniform first-amplitude estimate
in \(p\) or \(N_0\) is claimed.\\
\bottomrule
\end{longtable}

Every estimate stated for sufficiently small \(\eps\) means
\[
 \forall(p,a,b)\ \exists\eps_0,C_j,\ldots\
 \forall\,0<\eps<\eps_0\
 \forall\,N\text{ in the specified window}.
\]
Neither the point nor the derivative order is placed after the
\(\eps\)-quantifier. The \(N\)-window and \(\lambda\)-window splits
are made after differentiation; their moving endpoints are not
being differentiated.




\begin{thebibliography}{9}
\bibitem{Boukraa2008}
S.~Boukraa, A.~J.~Guttmann, S.~Hassani, I.~Jensen,
J.-M.~Maillard, B.~Nickel, and N.~Zenine,
Experimental mathematics on the magnetic susceptibility of the
square lattice Ising model,
\emph{J. Phys. A: Math. Theor.} \textbf{41} (2008), 455202.
\href{https://arxiv.org/abs/0808.0763}{arXiv:0808.0763v1}.
Physical normalization and complete angular representation:
equations (3), (5)--(9).

\bibitem{Orrick2001}
W.~P.~Orrick, B.~G.~Nickel, A.~J.~Guttmann, and J.~H.~H.~Perk,
The susceptibility of the square lattice Ising model: new developments,
\emph{J. Stat. Phys.} \textbf{102} (2001), 795--841.
\href{https://perk.okstate.edu/papers/canberra/chitex.pdf}{Author-hosted full text}.
Equations (1.9), (1.11), (3.13), (3.17)--(3.20), and (4.1)--(4.2).

\bibitem{TW2014}
C.~A.~Tracy and H.~Widom,
On the singularities in the susceptibility expansion for the
two-dimensional Ising model,
\emph{J. Stat. Phys.} \textbf{156} (2014), 1125--1135.
\href{https://arxiv.org/abs/1403.3966}{arXiv:1403.3966v1}.
The gradient argument is reconstructed here at the selected odd
orders; its footnote 5 is explicitly addressed.

\bibitem{Snapshot}
D.~Lin, \emph{A candidate natural-boundary proof for the bulk
square-lattice Ising susceptibility}, manuscript v0.1-rc6,
\texttt{paper/next/manuscript.tex}, in the user-supplied snapshot of
\href{https://github.com/LeoLam233/ising-bulk-natural-boundary}{ising-bulk-natural-boundary},
commit \texttt{2cf997226db314ad52fa40a2ecda5439542606bf},
17 October 2026. The current manuscript reuses and modifies its
explicit tail derivations, including all stated coordinate and
uniformity arguments, and does not assume its final theorem.
\end{thebibliography}

\end{document}
